% Outline: 6
\chapter{任务泛化与适应}
\label{chap:rl-transfer}

一辆已经学会绕过障碍的小车，换一个终点还需要从头试错吗？如果地图不变而地面变滑，旧的转向习惯还能直接使用吗？这两个问题看似都在问“能否迁移”，所需保留的知识却不同：更换目标通常不必重新认识动作的物理效果，改变移动规律却会使原先的到达概率失效。即使任务确实共享结构，智能体也未必知道新任务是哪一种，更未必获准通过试错弄清它。

第~\ref{chap:rl-learning}~章已经建立单个任务中的价值、策略与模型学习。本章增加任务之间的关系：先确定变化范围、可用信息与预算，再判断哪些知识可以保留、它们适用于哪些目标任务。训练阶段形成这些知识，部署阶段则依据描述、历史和更新权限形成具体策略。图~\ref{fig:rl-transfer-map}~把这几项职责分开；其中的验证反馈意味着失败可能要求修改共享假设，而不仅是增加训练次数。涉及历史充分性与固定日志时，分别调用第~\ref{chap:rl-information}~章的信念与覆盖条件；反馈含义仍由第~\ref{chap:rl-reward}~章的目标规定。

\begin{figure*}[t]
  \centering
  % Original chapter relationship diagram; final width 169 mm.
\begin{tikzpicture}[x=\linewidth,y=1mm,
  font=\figurefont\fontsize{8.5}{11}\selectfont,
  box/.style={draw=BookRule,fill=BookPaper,line width=.8pt,
    text width=28mm,minimum height=12mm,align=center,inner sep=2mm},
  flow/.style={-{Latex[length=2mm]},draw=BookInk,line width=.9pt},
  lab/.style={fill=white,inner sep=1mm,font=\figurefont\fontsize{8}{10}\selectfont}]
  \node[box] (conditions) at (.13,0) {任务、信息\\与预算};
  \node[box] (objects) at (.40,17) {共享结构\\与知识对象};
  \node[box] (range) at (.73,17) {目标任务\\适用范围};
  \node[box] (training) at (.40,-17) {共享知识训练};
  \node[box] (policy) at (.73,-17) {目标任务\\策略形成};
  \node[box,fill=BookTeal!8] (test) at (.73,-48) {验证与诊断};
  \draw[flow] (conditions.north) |- node[lab,pos=.7] {比较条件} (objects.west);
  \draw[flow] (objects) -- node[lab] {复用候选} (range);
  \draw[flow] (objects) -- node[lab] {学习对象} (training);
  \draw[flow] (range.east) -- (.97,17) -- node[lab,rotate=90] {测试范围}
    (.97,-65) -- (.73,-65) -- (test.south);
  \draw[flow] (range.south) -- node[lab,sloped,above] {训练覆盖} (training.north);
  \draw[flow] (training) -- node[lab] {知识供给} (policy);
  \draw[flow] (conditions.south) -- (.13,-33) --
    node[lab] {信息与更新权限} (.73,-33) -- (policy.south);
  \draw[flow] (policy.east) -- (.88,-17) |-
    node[lab,pos=.7,align=center] {行为\\与代价} (test.east);
  \draw[flow,dashed,draw=BookAmber] (test.west) -- (.02,-48)
    -- node[lab,rotate=90] {失败证据修正假设} (.02,32) -| (objects.north);
\end{tikzpicture}

  \caption{任务泛化与适应的论证关系。实线标明条件约束或知识供给，虚线表示验证后的修订反馈；不是所有方法都必须顺序执行的流程。部署中的记忆、参数更新和持续交互可以组合。图为本章关系的原创示意。}
  \label{fig:rl-transfer-map}
\end{figure*}

% Outline: 6.1
\section{任务泛化与适应问题}
\label{sec:rl-transfer-problem}

% Outline: 6.1.1
\subsection{任务要素的变化范围}
\label{sec:rl-transfer-components}

沿用前文宽4、高3、障碍位于$(2,2)$的导航地图，起点仍为$(1,1)$。基准任务以$(4,3)$为终点，合法意图方向以$0.8$概率执行，两侧各以$0.1$概率滑移；进入终点得1，其他非终止步得$-0.04$，吸收后奖励为0，折扣为$\gamma=0.9$。这些是统一的教学设定，不是实测数据。现在分别改变终点、移动噪声、观测符号或动作编码，检查哪些数学分量随之改变。旧策略可能依次错在“去哪里”“动作会到哪里”“眼前是什么”和“命令代表什么”，而不只是遇到了不同的数据文件。

\begin{definition}[任务族]
  \label{def:rl-transfer-family}
  用任务参数$c\in\mathcal C$索引环境
  \[
    \mathcal M_c=(\mathcal S_c,\mathcal A_c,P_c,R_c,O_c,\rho_{0,c},\gamma).
  \]
  其中$O_c$为观测核，$R_c(s,a)$为动作时刻奖励的条件期望。任务族是这些环境及其任务信息接口的集合。比较时须声明哪些分量随$c$改变，哪些共同固定；有终止事件或有限时域时，还须规定终止机制与时域。
\end{definition}

同一状态集合上的两个奖励函数已经定义了不同任务；状态和动作的维数相同，并不要求转移、奖励或语义相同。若要根据随机奖励和终止反馈辨识任务，还须给出这些反馈与后继状态的联合条件分布：均值$R_c$本身不能确定反馈似然。表~\ref{tab:rl-transfer-changes}~列出四种单因素变体，其中可直接观察的信息由实验协议提供，不是把智能体训练中见过的信息一概当作已知。

\begin{table*}[t]
  \centering\small
  \begin{tabularx}{\linewidth}{@{}>{\raggedright\arraybackslash}p{23mm}*{3}{>{\raggedright\arraybackslash}X}@{}}
    \toprule
    变化分量 & 保持不变的对象 & 可直接观察的信息 & 需要重算或重新学习的对象\\
    \midrule
    终点与奖励 & 地图、物理移动规律 & 可能提供目标坐标；也可能只有到达反馈 & 新目标价值、路线选择；终止规则改变时还包括有效转移\\
    移动噪声 & 目标含义、地图、命令语义 & 位置与移动结果；风参数未必给出 & 转移概率、长期后果、控制策略\\
    观测符号 & 潜在地图与物理动作 & 新符号图像；对应表可能未知 & 感知对齐，必要时重学历史表示\\
    左右命令互换 & 物理空间、目标、观测 & 命令标签与执行结果 & 动作映射、合法命令集与技能接口\\
    \bottomrule
  \end{tabularx}
  \caption{同一导航地图中的任务变化。每行单独改变一个因素，数字与任务均为教学构造。更换吸收终点会改变完整MDP的转移核，不能只因物理移动规则相同就声称共同动力学。}
  \label{tab:rl-transfer-changes}
\end{table*}

这里有一个容易遗漏的细节：如果“终点”意味着进入后吸收，那么换终点不仅换奖励，也改变哪些状态可以继续移动。需要严格隔离奖励变化时，本章会显式改用共同终止规则：例如所有任务在相同时间结束，或使用固定的吸收出口、只改变沿途地标的奖励。前一种有限时域设定还须把剩余时间纳入状态，才能用同一个转移核表示终止机制。这些是为检验复用条件而定义的独立变体，不是基准网格的另一种数值解。

还要区分两个时间尺度。在一次适应试验中，未知风向可以固定不变，只是智能体逐渐认识它；也可以每个回合都改变。前者的证据指向同一估计对象，后者的旧证据可能已经过时。固定任务辨识在第~\ref{sec:rl-transfer-identification}~节展开，变化任务的跟踪在第~\ref{sec:rl-transfer-tracking}~节另行定义。

% Outline: 6.1.2
\subsection{目标任务信息的可用条件}
\label{sec:rl-transfer-information}

终点标在地图上时，小车可以立即规划；终点隐藏、抵达才发奖励时，它必须从反馈中排除候选位置；风向隐藏时，即使终点已知，也可能要试一次移动才能分辨转移规律。目标是否变化与任务身份是否可见是两个问题。已知新目标不等于已知新动力学，未知任务也不等于没有任何先验。

\term{显式任务描述}（explicit task description）是执行前提供的目标坐标、参数或语言指令，其作用是直接约束任务含义。\term{历史任务信息}是由已记录的观测、动作、奖励与终止标记推断的任务信息。\term{主动辨识}（active identification）则通过选择动作获得更有区分力的反馈；它改变了未来可用的信息，而不仅是对已有历史多计算几遍。语言描述只有在其与奖励和控制接口的对应关系成立时才有用，不能把一句自然语言自动当作完整MDP。

设训练任务用编号1和2分别代表东、西两个终点。把编号交换、保留所有轨迹的物理内容，再相应交换训练标签，并不会改变问题难度。对于从未出现的编号3，训练记录没有告诉模型它是北端终点还是强风任务。编号能索引记住的任务，却没有坐标那样的几何语义。新编号上的零样本迁移必须另有描述映射，不能依靠编号的大小顺序。

历史编码器输出的向量$\bm z_t$同样只是一个计算结果。若它经过奖励预测或策略损失训练，可以包含有用的任务信息，但不自动满足后验概率的归一化与校准含义。更根本地，若所有获准动作在两个任务中都产生相同的反馈分布，再长的历史和再大的模型也不能凭空区分任务。第~\ref{sec:rl-transfer-belief}~节会把这种信息缺口与推断近似误差分开。

% Outline: 6.1.3
\subsection{零样本、少样本与持续适应预算}
\label{sec:rl-transfer-budget}

“不更新网络”并不是完整的预算说明。一个网络可以保持权重不变，却在上下文中读过大量目标任务反馈；另一个网络可能在固定日志上更新参数，却没有权限增加一次真实交互。为比较两者，需要分别记录数据、可变状态和计算。

本章将\term{零样本执行}（zero-shot execution）限定为不使用目标任务反馈进行额外学习、记忆适应或模型选择的执行。它仍可读取给定任务描述和正常控制所需的观测。部分可观测任务中的状态滤波应预先声明：在已给定任务模型下估计当前位置，属于控制所需的状态估计；利用偏移反馈推断未知风向，则改变了任务信息，属于适应，即使只发生在一个回合内、网络权重始终不变。\term{少样本适应}（few-shot adaptation）使用预定数量的目标经验形成策略；经验可以来自给定日志，也可以来自获准交互。\term{持续适应}（continual adaptation）允许在部署过程中反复更新任务信息或策略，评价必须计入这一过程中的表现。

表~\ref{tab:rl-transfer-budget}~用教学协议固定这些差别。$B$表示适应交互步数，$L$表示最多保留的历史长度，$K$表示梯度步数；实际论文还应报告各回合长度，不能仅以“若干回合”比较不同任务。

\begin{table*}[t]
  \centering\small
  \begin{tabularx}{\linewidth}{@{}>{\raggedright\arraybackslash}p{19mm}*{6}{>{\raggedright\arraybackslash}X}@{}}
    \toprule
    协议 & 目标交互 & 可见反馈 & 上下文 & 参数更新 & 额外计算 & 选模查询\\
    \midrule
    零样本 & 无适应步；执行另计 & 描述、执行观测 & 不累计任务学习信息 & 0 & 固定推理预算 & 0\\
    固定日志适应 & 新交互0 & 日志给定字段 & 至多$L$步 & 至多$K$步 & 推断与反传均计 & 预先固定\\
    少量在线适应 & 总计$B$步 & 约定奖励及终止 & 至多$L$步 & 至多$K$步 & 探测、规划均计 & 纳入$B$\\
    持续在线适应 & 全程累计 & 按到达时刻可见 & 记忆重置规则固定 & 频率与总量均报 & 每步和累计成本 & 不用未来成绩选模\\
    \bottomrule
  \end{tabularx}
  \caption{目标任务预算的不同分量。各列共同定义协议；“固定日志”与“固定参数”不是同义词。本表为评价口径，不是算法性能数据。}
  \label{tab:rl-transfer-budget}
\end{table*}

假设每回合恰有20步，共运行10回合，这是教学预算而非实测配置。只报告第1回合、按固定控制规则响应观测且不辨识任务或更新策略，回答的是初始执行能力；允许在10回合内学习并报告全部200步回报，回答的是学习期间能否付得起试错成本；先适应10回合，再停止任务学习、另测5回合，回答的是200步适应后的控制能力，但总交互已是300步。三种分数即使相等，也没有相同含义。

试探风向、从失败位置恢复、查询检查点成绩都消耗目标任务资源。选择“最好的第7次更新”需要看目标任务结果，不能只报告最终执行的步数。预算的完整性决定了“快”的比较对象。

% Outline: 6.2
\section{可复用决策知识的确定}
\label{sec:rl-transfer-knowledge}

% Outline: 6.2.1
\subsection{由任务要素异同判断共享结构}
\label{sec:rl-transfer-sharing}

只换导航目标时，旧地图仍能回答向右会走到哪里；把左右命令互换时，地图虽然相同，原策略发出的“右”却执行向左。直接复制参数是否可行，取决于参数所计算的对象，而不取决于网络能否接受同样形状的张量。图~\ref{fig:rl-transfer-components}~用逐分量矩阵显示这种区别。

\begin{figure}[htbp]
  \centering
  % Original controlled task matrix; symbols redundantly encode shared cells.
\begin{tikzpicture}[x=\linewidth,y=1mm,
  font=\figurefont\fontsize{8.5}{11}\selectfont,
  cell/.style={draw=BookRule,line width=.7pt,minimum width=17mm,
    minimum height=8mm,inner sep=1mm,fill=BookPaper}]
  \foreach \x/\name in {.35/物理转移,.54/计分规则,.73/观测,.92/动作语义}
    \node at (\x,8) {\name};
  \foreach \y/\name in {0/基准,-10/换目标,-20/换噪声,-30/换观测,-40/换命令}
    \node[anchor=east] at (.23,\y) {\name};
  \foreach \y in {0,-10,-20,-30,-40} {
    \node[cell] at (.35,\y) {$P$};
    \node[cell] at (.54,\y) {$f_r$};
    \node[cell] at (.73,\y) {$O$};
    \node[cell] at (.92,\y) {$A$};
  }
  \foreach \x/\y/\value in {.54/-10/{f_{r,g}},.35/-20/P_\eta,.73/-30/O_c,.92/-40/A_c} {
    \node[cell,fill=BookTeal!10] at (\x,\y) {$\value$};
    \draw[BookTeal,line width=.8pt] (\x-.07,\y-3) -- (\x-.05,\y-1);
  }
  \node[anchor=north west,align=left] at (.02,-49)
    {共同终止规则；$P$按物理动作索引。\\
     $A_c$表示命令到物理动作的映射。};
\end{tikzpicture}

  \caption{导航任务族的要素矩阵。相同符号与底纹表示同一分量，斜线及新符号表示变化；“换目标”行使用共同终止规则。$f_r(s,a,s')$表示给定转移后的计分规则，换目标时改为$f_{r,g}$；换噪声保留$f_r$，却可能改变均值奖励$R_c(s,a)$。$P$按物理动作索引，命令交换由$A_c$表示；若按命令标签索引，转移核也须相应变换。矩阵为控制变量设计，不表示从数据证明了共享关系。}
  \label{fig:rl-transfer-components}
\end{figure}

\term{共享结构}（shared structure）是关于任务对象的可检查假设。例如，在统一状态动作语义下$P_c=P$；奖励都由同一组特征线性组合；观测可以映射到同一控制充分状态；或技能在各任务中具有同样的启动集与终止含义。这些假设强弱不同，不能统称为“任务相似”后省略。

实验者编写的任务生成器可以保证某个模块不变，此时结构是已知的。若从多任务轨迹估计共同转移，再通过留出数据检验，则结构是估计结果，带有覆盖与误差条件。仅因两条学习曲线相似而猜测可以共享，只是待检验假设。控制变量比较应在相同状态、相同物理动作及相同反馈定义上检查差别；否则观测编码错误可能被误诊为动力学变化。

相似特征与相似回报都只能提供线索。一个编码器可能把两个外观相近、却需要相反动作的状态合并；两项任务也可能恰好在当前策略访问的区域一致，而在未访问区域完全不同。共同结构的声明必须带上适用状态动作范围。

% Outline: 6.2.2
\subsection{由共享结构确定可复用知识}
\label{sec:rl-transfer-objects}

共享关系确定后，才有理由选择保留什么。转移模型描述给定状态和动作后的变化规律，因此动作的物理效果不变时，它可能继续适用。价值函数评价的是在特定动力学与奖励下遵循某份策略的长期回报；只更换奖励，也可能需要重算价值。策略则根据可用信息选择动作：即使它在新任务中仍能正常运行，也要检查这些动作能否实现新目标。判断复用范围，需要逐项核对所保存的预测或决策规则依赖哪些任务条件。

例如，把终点从$(1,3)$换到$(4,1)$，可保留物理地图与动作效果，但到旧终点的距离和价值不能直接解释新任务。把滑移概率由基准的$0.2$改成教学假设$0.4$，目标坐标仍有相同含义，但旧到达概率必须重新估计。表~\ref{tab:rl-transfer-objects}~将复用对象与所需条件成对列出。

\begin{table*}[t]
  \centering\small
  \begin{tabularx}{\linewidth}{@{}>{\raggedright\arraybackslash}p{21mm}*{3}{>{\raggedright\arraybackslash}X}@{}}
    \toprule
    复用对象 & 共享条件 & 目标任务仍需估计的量 & 典型失效\\
    \midrule
    状态表示 & 保留目标控制所需信息；观测可对齐 & 任务特定变量、价值或策略 & 丢掉新目标需要的线索\\
    环境模型 & 相同转移机制或已知模块关系 & 风向、摩擦等特定参数 & 未建模机制改变\\
    奖励特征 & 特征语义不变，新奖励可表示 & 新权重$\bm w_c$ & 新目标不在特征类中\\
    后继表示与特征 & 固定策略、动力学和折扣 & 新奖励权重；必要时补新策略 & 转移变化或策略库覆盖差\\
    价值函数 & 原价值对应目标奖励、转移与策略 & 不满足条件时需重估价值 & 旧数值仍可算但含义已错\\
    策略初始化 & 接口兼容，少量更新有合适方向 & 目标任务参数修正 & 小数据噪声导致负适应\\
    技能 & 启动、输入、动作效果和终止兼容 & 高层选择、组合或内部适应 & 能启动却不能达到所需效果\\
    \bottomrule
  \end{tabularx}
  \caption{可复用知识及其条件。共享条件说明对象何时有意义，不承诺有限数据训练能够满足这些条件。来源为本章定义与后文推导的归纳。}
  \label{tab:rl-transfer-objects}
\end{table*}

技能是具有持续时间的行为，而不是一个固定长度动作数组。即使“绕过障碍”技能仍能运行，新任务也可能要求中途停在障碍旁；此时技能的终止语义不合适。检查启动集、输入、控制效果与终止条件，比仅运行一次确认“没有报错”更接近决策层面的兼容性。时间抽象的单任务定义见第~\ref{chap:rl-learning}~章，本章关心这些条件如何跨任务保持。

% Outline: 6.3
\section{可复用决策知识的适用任务}
\label{sec:rl-transfer-applicability}

同一个复用对象可以在不同目标任务上成立或失效。本节固定候选对象，分别检验任务生成关系、奖励、动力学以及观测动作接口；这些是并列的检查视角，可以同时作用于一个任务。

% Outline: 6.3.1
\subsection{训练任务与目标任务的生成关系}
\label{sec:rl-transfer-distributions}

从一张已知地图随机抽取更多轨迹，与从多种目标和风参数中抽取新任务，是两层不同的采样。用$p_{\text{训练}}(c)$表示训练任务总体分布，$p_{\text{目标}}(c)$表示所声明部署范围的任务分布。给定任务$c$后，轨迹还服从由$P_c$、初始分布和收集策略共同决定的分布$p_c^\mu(\tau)$。因此训练平均至少包含
\[
  \mathbb E_{c\sim p_{\text{训练}}}
  \mathbb E_{\tau\sim p_c^\mu}[\ell(c,\tau)].
\]
外层变动任务，内层变动同任务中的经历；大量内层样本不能替代外层任务多样性。

有限任务清单$\{c_1,\ldots,c_n\}$只给出了实际训练对象。任务生成器是可继续采样新对象的程序，其采样规则定义一个分布。理论中的总体分布则可能是对部署环境的假设，并不一定有可运行生成器。使用经验分布$n^{-1}\sum_i\delta_{c_i}$时应明确它只在已见任务上有质量，不能把这个支持集与假设总体的支持混用。

\term{上下文MDP}（contextual MDP）用$c$控制一个MDP的分量；这里的“上下文”指任务条件，未必是网络的历史窗口。物理状态与$c$都可见时，策略可以写为$\pi(a\mid s,c)$。物理状态可见、$c$隐藏但在适应期间固定时，$s_t$加上任务后验$b_t(c)$形成\term{贝叶斯自适应MDP}（Bayes-adaptive MDP）的状态：动作既影响位置，也影响对任务的认识。状态也隐藏时则需联合信念，具体更新见第~\ref{sec:rl-transfer-belief}~节。

图~\ref{fig:rl-transfer-task-plane}~给出一组教学任务划分。横轴为地图上边界目标$(x,3)$的横坐标，纵轴为总滑移概率$\eta$。训练总体仅允许$x\in\{1,2,3\}$，其中$x=3$只与$\eta\leq0.2$组合；$x=1,2$允许$0\leq\eta\leq0.4$。抽到一个未见的$(1,0.15)$仍属于同分布新任务；$(2,0.2)$可检验已见噪声值之间的插值；$(3,0.3)$的两个因素各自出现过，却在联合支持之外；$(4,0.45)$则含未见因素范围。

\begin{figure}[htbp]
  \centering
  % Teaching assumptions only; x is discrete goal coordinate, eta is slip probability.
% Population support: x=1,2 with eta in [0,.4]; x=3 with eta in [0,.2].
\begin{tikzpicture}
  \begin{axis}[width=.98\linewidth,height=65mm,
    font=\figurefont\fontsize{8.5}{11}\selectfont,
    axis lines=left,xmin=.6,xmax=4.4,ymin=0,ymax=.52,
    xtick={1,2,3,4},ytick={0,.1,.2,.3,.4,.5},
    xlabel={目标$(x,3)$的横坐标$x$},ylabel={总滑移概率$\eta$},
    tick label style={font=\figurefont\fontsize{8}{10}\selectfont},
    clip=false]
    \path[fill=BookTeal!10] (axis cs:.88,0) rectangle (axis cs:1.12,.4);
    \path[fill=BookTeal!10] (axis cs:1.88,0) rectangle (axis cs:2.12,.4);
    \path[fill=BookTeal!10] (axis cs:2.88,0) rectangle (axis cs:3.12,.2);
    \addplot[only marks,mark=*,mark size=2pt,BookInk]
      coordinates {(1,.1)(1,.3)(2,.1)(2,.3)(3,.1)};
    \addplot[only marks,mark=o,mark size=3pt,BookBlue,thick]
      coordinates {(1,.15)};
    \addplot[only marks,mark=square,mark size=3pt,BookTeal,thick]
      coordinates {(2,.2)};
    \addplot[only marks,mark=triangle,mark size=3pt,BookAmber,thick]
      coordinates {(3,.3)};
    \addplot[only marks,mark=x,mark size=3pt,BookInk,thick]
      coordinates {(4,.45)};
    \node[anchor=west,font=\figurefont\fontsize{8}{10}\selectfont]
      at (axis cs:1.13,.15) {同分布};
    \node[anchor=west,font=\figurefont\fontsize{8}{10}\selectfont]
      at (axis cs:2.13,.2) {插值};
    \node[anchor=south,font=\figurefont\fontsize{8}{10}\selectfont]
      at (axis cs:3,.32) {未见组合};
    \node[anchor=east,font=\figurefont\fontsize{8}{10}\selectfont]
      at (axis cs:3.9,.47) {新因素范围};
    \node[anchor=west,font=\figurefont\fontsize{8}{10}\selectfont]
      at (axis cs:.8,.49) {实点：训练样本};
  \end{axis}
\end{tikzpicture}

  \caption{目标位置与移动噪声的教学任务平面。竖向浅带表示离散目标坐标下的总体噪声支持，实点为有限训练样本，其他点形区分留出任务类型。插值只沿噪声轴解释，不把离散网格目标当成连续位置；所有点与支持范围均为教学设定。}
  \label{fig:rl-transfer-task-plane}
\end{figure}

这里“未见”相对于有限样本，“支持外”相对于声明的总体，二者不是同义词。若原总体本来就独立抽取目标与噪声，$(3,0.3)$可能只是在训练清单中缺席，而不在总体支持之外。组合测试必须说明所依据的是边缘支持、联合支持还是经验清单。

% Outline: 6.3.2
\subsection{目标与奖励变化下的可复用知识}
\label{sec:rl-transfer-reward-change}

保持地图与移动规律不变，把目标从左上角移到右下角，重新规划不必从头估计每条物理转移。但旧价值把未来访问按旧目标计分，不能只改名字继续使用。若希望不重新展开每条轨迹就评价新奖励，需要保留尚未加权的长期后果。

设所有任务有共同$\mathcal S,\mathcal A,P,\gamma$以及一致的观测、动作和终止语义。将转移条件下的期望奖励写成
\begin{equation}
  R_c(s,a,s')=\langle\bm\phi(s,a,s'),\bm w_c\rangle,
  \qquad \bm\phi,\bm w_c\in\mathbb R^d.
  \label{eq:rl-transfer-linear-reward}
\end{equation}
这里$R_c(s,a)=\mathbb E_{s'\sim P(\cdot\mid s,a)}R_c(s,a,s')$，没有把奖励样本$r_t$当成已知期望。共同的$\bm\phi$记录“经过了哪些地标、付出了多少步成本”等事件，$\bm w_c$规定新任务怎样评价这些事件。

在固定策略下，保存未来各特征的期望折扣总量，就能用新权重重新计分。这一对象称为\term{后继特征}（successor features, SF）。对于教学路径，若前两步分别出现$\bm\phi_0=(1,0)^\top$、$\bm\phi_1=(0,1)^\top$，之后特征为0，那么$\gamma=0.9$时后继特征为$(1,0.9)^\top$；权重$(2,-1)^\top$给出价值$1.1$，换为$(0,3)^\top$则给出$2.7$。变化的是评价标准，而不是路径分布。

共同终止规则在这里不可省略。若新目标一到达就吸收，路径分布也随之改变，旧后继特征不能自动表示新任务。可改用固定吸收出口、仅改变经过地标的奖励，或重新学习包含新终止机制的后继特征。第~\ref{sec:rl-transfer-value-training}~节会完整推导价值分解，并讨论如何从多条已有策略中逐状态选择动作。原始后继特征框架正是针对共同动力学下的奖励迁移\scite{rl6-barreto2017-successor}

即使新奖励可表示，策略库也可能没有掌握通往新目标的行为。保存了正确的后果表示，只能评价已有行为及其可用组合，不等于已经拥有新任务的最优行为。

% Outline: 6.3.3
\subsection{动力学变化下的可复用知识}
\label{sec:rl-transfer-dynamics-change}

固定终点，只把可靠移动改成受风影响的随机移动，进入终点得1的含义没有变，但选择一条窄路的长期后果会变。动力学复用应区分模型结构与模型参数。若风参数$\eta$已提供，而且模型族$P_\eta$正确，可以代入新参数重新规划；若$\eta$未知但只是一维参数，则可能只需辨识它，而不必独立估计每个状态动作的全部后继概率。

例如，在没有边界或障碍干扰的指定测试位置，教学模型令移动偏离意图的指示量为$X_j\sim\operatorname{Bernoulli}(\eta)$。独立重置采样$n$次时，估计$\widehat\eta=n^{-1}\sum_jX_j$的方差为$\eta(1-\eta)/n$。将这一参数用于全图，需要额外假设风机制处处相同，且已正确处理撞墙后留在原位的合并概率。若只观察墙边位置，若干不同滑移方向可能产生相同后继状态，不能再把每条记录直接当作可观测的$X_j$。

共享模块可以表达“几何碰撞规则相同，风模块变化”；平滑模型可以表达“小幅改变$\eta$只小幅改变转移”。前者减少重复估计的对象，后者约束插值误差，两者都不是从网络参数共享自动推出的事实。一般有限表格转移要估计约$|\mathcal S||\mathcal A|(|\mathcal S|-1)$个自由概率，单个风参数只估计一个量，但这个维数比较不包含辨识它所需的访问成本。

反例更直接：在同一可见岔路，两个动作分别通向奖励1和0的区域；第二个任务将两个动作的物理效果互换，图像与奖励标记完全不变。旧策略仍选原来的动作，回报由1变为0。旧后继特征预测它会访问奖励区域，因而同样失效；仅修改奖励权重不能把错误转移改对。这个单步例是从导航提取的教学简化，取消滑移与步成本。

模型误差还会传播到控制。若两个模型共享条件期望奖励$R(s,a)$，满足$|R|\leq R_{\max}$，且转移满足$\sup_{s,a}\lVert P_c-P_{c'}\rVert_1\leq\delta$，固定策略的Bellman关系给出
\[
 \lVert V_c^\pi-V_{c'}^\pi\rVert_\infty
 \leq \gamma\lVert V_c^\pi-V_{c'}^\pi\rVert_\infty
      +\gamma\delta\,\frac{R_{\max}}{1-\gamma}.
\]
移项得到上界$\gamma\delta R_{\max}/(1-\gamma)^2$。若只是“进入终点得1”的规则相同，改变到达概率仍可能改变$R_c(s,a)$；两者仍满足同一$R_{\max}$上界时，还要在价值界中加入$\lVert R_c-R_{c'}\rVert_\infty/(1-\gamma)$，不能直接套用共同均值奖励的特例。这说明长期后果可能放大小的单步差异；它要求全相关状态动作的一致误差，训练参数范围内的平均预测误差不能代替该条件，更不能替范围外参数担保。

% Outline: 6.3.4
\subsection{观测与动作接口变化下的可复用知识}
\label{sec:rl-transfer-interface}

把地图颜色换掉，并不一定改变小车的物理位置；交换左右动作标签，则会改变同一输出命令的控制效果。对一段“右、右、上”的命令序列，单纯换色后物理轨迹可以完全不变，交换命令后却可能向相反方向移动，甚至使命令不再合法。因此输入与输出应分别对齐。

设新观测编码为$e_c(o)$，它映射到原策略使用的决策表示；动作映射$m_c(a)$把原语义动作变为新环境命令。若潜在状态之间有对应$f_c$，精确复用的一个充分条件是奖励保持且
\[
 P_c\!\left(f_c^{-1}(B)\mid s,m_c(a)\right)
 =P_0\!\left(B\mid f_c(s),a\right)
\]
对所有相关状态、合法动作与后继集合$B$成立，同时$e_c$保留决策所需信息、$m_c$保持合法性。这是控制效果的对应，不只是颜色或向量之间的相似。图~\ref{fig:rl-transfer-interface}~分别标出编码、决策、映射与执行。

\begin{figure*}[t]
  \centering
  % Original interface diagram; final width 169 mm.
\begin{tikzpicture}[x=\linewidth,y=1mm,
  font=\figurefont\fontsize{8.5}{11}\selectfont,
  box/.style={draw=BookRule,fill=BookPaper,line width=.8pt,
    text width=22mm,minimum height=12mm,align=center,inner sep=1.5mm},
  flow/.style={-{Latex[length=2mm]},draw=BookInk,line width=.9pt}]
  \node[box] (obs) at (.08,0) {新观测\\$o_t$};
  \node[box,fill=BookTeal!8] (enc) at (.29,0) {对齐表示\\$e_c(o_t)$};
  \node[box] (policy) at (.50,0) {原策略\\$\pi_0$};
  \node[box,fill=BookTeal!8] (act) at (.71,0) {动作映射\\$m_c$};
  \node[box] (env) at (.92,0) {新环境\\$P_c$};
  \draw[flow] (obs) -- node[above] {编码} (enc);
  \draw[flow] (enc) -- node[above] {决策} (policy);
  \draw[flow] (policy) -- node[above] {映射} (act);
  \draw[flow] (act) -- node[above] {执行} (env);
  \draw[flow] (env.south) -- ++(0,-13) -|
    node[pos=.25,below] {新观测与反馈} (obs.south);
  \node[anchor=north,align=center,text=BookMuted] at (.5,-23)
    {输入：保留控制所需信息\qquad 输出：保留动作效果、合法性与技能终止语义};
\end{tikzpicture}

  \caption{观测和动作接口的迁移。输入编码需要保留控制信息，输出映射需要保留物理效果与合法性；任何一侧缺少对应关系，原策略都不能原样解释新接口。图为上述充分条件的原创结构示意。}
  \label{fig:rl-transfer-interface}
\end{figure*}

对于技能，还要把启动状态映射到原启动集，把“到达出口”映射到相同的终止事件。只映射原始动作，而忽略技能结束条件，可能导致高层策略等待一个永远不会到来的反馈。

有限日志未必能唯一确定这些映射。若日志只记录“上”和“下”，左右动作交换前后与全部记录相容；图像重建也无法辨认未试过命令的含义。需要增加有区分力的动作记录，或者提供可信的接口说明。控制充分表示与数据可识别性的区别见第~\ref{chap:rl-information}~章：同尺寸接口、感知相似度和重建成功都不够。

% Outline: 6.3.5
\subsection{组合与支持外任务所需的额外假设}
\label{sec:rl-transfer-composition}

训练时见过“远终点、无风”和“近终点、有风”，并不意味着已经见过“远终点、有风”。若目标模块只改变奖励、风模块只改变转移，且二者接口独立，模型可以组合这两种知识；若有一个仅在远终点且有风时出现的障碍，原来的两项任务却完全不能揭示它。

形式上，令$c=(g,\eta)$。边缘支持分别为$\operatorname{supp}(p_g)$和$\operatorname{supp}(p_\eta)$，联合支持则为$\operatorname{supp}(p_{g,\eta})$。一般只有
\[
 \operatorname{supp}(p_{g,\eta})
 \subseteq \operatorname{supp}(p_g)\times\operatorname{supp}(p_\eta),
\]
不保证相等。\term{组合泛化}（compositional generalization）要求把已学因素按未见组合使用；成立依据可能是$P_c=P_\eta$、$R_c=R_g$等因子分解，以及组合后仍成立的观测、动作与终止接口。这些条件必须单独验证。

可以构造两个生成器：它们在已见组合上给出完全相同的地图，只在“远且有风”时分别保持通路或封闭通路。任何只用现有训练组合的算法面对两者都接收相同数据，因而不能仅凭这些数据确定新组合是否可通行。增加同一组合的轨迹数量无法消除这个歧义。

从未见过的新动作效果则更强：即使所有旧因素能正确组合，也没有定义它怎样作用于环境。支持外外推需要可识别的机制、适用范围明确的平滑性或结构不变性。测试应分别留出组合与因素范围；“训练了很多任务”不能代替对这种关系的证据。

% Outline: 6.4
\section{可复用决策知识的训练}
\label{sec:rl-transfer-training}

训练任务和经验已经给定后，需要选择它们进入哪个预测或决策目标。任务随机化、条件化、蒸馏和目标重标记是训练手段；表示、模型、价值和策略才是被训练的对象。同一种手段可以服务多个对象，其有效性仍受前两节的共享条件与任务支持约束。

% Outline: 6.4.1
\subsection{训练任务与经验的构造}
\label{sec:rl-transfer-data}

参数随机化可以每次抽取目标坐标和风强度；程序化关卡还可以抽取地图，但必须检查连通性、合法动作和终止规则。无论任务生成多复杂，得到经验仍有两个步骤：按$q(c)$抽取任务，再由该任务中的收集策略$\mu_c$产生轨迹。折扣情形下，若用归一化占用度量抽样，经验权重为$q(c)d_c^{\mu_c}(s,a)$。某任务被频繁抽到，不代表其关键窄路也被充分访问。

任务采样权重还会改变优化目标。若希望优化$p_{\text{训练}}$下的平均量，却用$q$抽任务，需要$p_{\text{训练}}$相对于$q$绝对连续，并使用相应概率或密度比$p_{\text{训练}}(c)/q(c)$。离散情形下，这要求每个$p_{\text{训练}}(c)>0$的任务都有$q(c)>0$；从未采样的任务不能靠重加权补回。权重过大会增加方差。课程学习选择当前较难的任务时，也应报告实际任务权重，而不是把采样改变隐藏在“多训练一些”之中。

有些失败轨迹无需重新交互，就能回答另一个目标的问题。\term{事后经验回放}（hindsight experience replay, HER）把轨迹实际达到的位置当作替代目标，重算该目标下的反馈，再交给离策略学习器。它重新解释了已经发生的转移，没有制造另一条物理轨迹\scite{rl6-andrychowicz2017-her}

例如，教学轨迹从$(1,1)$出发，依次到达$(2,1)$、$(3,1)$和$(4,1)$，没有到达原目标$(4,3)$。用新目标$g'=(3,1)$，设到达目标得1、之前每步得$-0.04$，则前两条转移的奖励由$(-0.04,-0.04)$改为$(-0.04,1)$。若新任务到达目标立即结束，第二条转移的终止标记必须同时改为1，后续从$(3,1)$继续移动的记录不能接在这条新回合后作为有效经验。图~\ref{fig:rl-transfer-relabel}~区分真实采集与离线重标记。

\begin{figure}[htbp]
  \centering
  % Original HER data-flow diagram. No relabeling edge creates real interactions.
\begin{tikzpicture}[x=\linewidth,y=1mm,
 font=\figurefont\fontsize{8.5}{11}\selectfont,
 box/.style={draw=BookRule,fill=BookPaper,line width=.8pt,
   text width=23mm,minimum height=11mm,align=center,inner sep=1.5mm},
 flow/.style={-{Latex[length=2mm]},draw=BookInk,line width=.9pt}]
 \node[box] (task) at (.14,0) {任务生成\\目标$g$};
 \node[box] (env) at (.50,0) {真实交互\\$\mu_g,P$};
 \node[box] (trajectory) at (.86,0) {轨迹\\$(s_t,a_t,s_{t+1})$};
 \node[box,fill=BookTeal!8] (relabel) at (.50,-26) {目标重标记\\$g'=(3,1)$};
 \node[box] (samples) at (.14,-26) {训练样本\\重算$r_t,\delta_t$};
 \draw[flow] (task) -- node[above] {采样} (env);
 \draw[flow] (env) -- node[above] {记录} (trajectory);
 \draw[flow] (trajectory.south) |- node[pos=.7,above] {已达位置} (relabel.east);
 \draw[flow] (relabel) -- (samples);
 \node[anchor=north,align=center,text=BookMuted] at (.5,-38)
   {上行：取得新经验\qquad 下行：重新解释已有经验};
\end{tikzpicture}

  \caption{任务生成与目标重标记的两条数据路径。上行产生真实轨迹，下行重算已有轨迹的目标、奖励和终止标记；重标记箭头不回连环境，因而不增加真实状态动作覆盖。原理据HER，具体导航路径为本章教学构造。}
  \label{fig:rl-transfer-relabel}
\end{figure}

重标记的语义有效性至少要求：目标不改变非终止物理转移，奖励及目标相关终止信号可按记录重算，动作在新解释中仍合法。若开门目标同时改变门是否锁住，就不能把旧转移当作新目标样本。即使语义有效，随机环境中的选择偏差仍是另一个问题。

设教学单步任务的同一个动作以各$1/2$概率落在左、右地标，且目标不影响转移。真实左目标下成功概率为$1/2$；如果只把“实际落在左边”的记录重标成左目标，再用这些记录估计该动作的左目标成功率，就得到1。每条保留记录都可能真实发生，但条件选择删除了失败分支。确定性动力学可避免这种特定的随机结果选择问题，随机环境则需检查重标记分布或采用适当校正；“目标不改变转移”本身不能证明估计无偏。

重标记$m$个目标会把每条原始记录扩成至多$m$条训练记录，奖励重算约需$O(m)$次调用。增加的是训练计算与目标覆盖，不是$m$倍独立环境证据。

% Outline: 6.4.2
\subsection{可复用状态表示的训练}
\label{sec:rl-transfer-representation-training}

不同终点任务都需要辨认障碍位置，但不应把目标位置压成同一个常量。所谓任务不变表示，应保留共同的控制要素，同时为会改变决策的任务变量留下入口。若表示把“同一地图”误解为“同一最优动作”，共享就会制造冲突。

共享编码器写作$\bm z=e_{\bm\theta}(o)$，策略或价值另接任务描述$c$；任务条件表示写作$\bm z=e_{\bm\theta}(o,c)$，在编码时就解释目标；模块化表示可以写作$\bm z=(e_{\bm\theta}(o),e_{\bm\omega_c}(o,c))$，将共同几何与特定传感或动力学信息分开。三种设计都可输入$Q(\bm z,a,c)$、$\pi(a\mid\bm z,c)$或潜在转移模型，表~\ref{tab:rl-transfer-representations}~固定这些用途比较。

\begin{table*}[t]
  \centering\small
  \begin{tabularx}{\linewidth}{@{}>{\raggedright\arraybackslash}p{18mm}*{5}{>{\raggedright\arraybackslash}X}@{}}
    \toprule
    设计 & 共享部件 & 任务特定部件 & 训练信号 & 任务权重 & 适用条件\\
    \midrule
    共享编码器 & 感知与几何表示 & 条件头或独立头 & TD误差、策略梯度、预测 & 显式$q(c)$ & 编码未丢目标相关信息\\
    条件表示 & 编码规则和参数 & 描述$c$或历史变量 & 条件价值、控制与预测 & 相同总量下检查各任务 & 描述语义及条件支持兼容\\
    模块化表示 & 共同子模块及组合器 & 风、传感器等模块 & 模块预测与端到端控制 & 可按模块或任务平衡 & 分解与组合接口成立\\
    \bottomrule
  \end{tabularx}
  \caption{三种表示设计。训练信号可以组合，但“参数共享”只描述实现，不能替代控制充分性证明。各列为本章机制归纳。}
  \label{tab:rl-transfer-representations}
\end{table*}

例如，以共享编码器输入一个评论家，一条经验给出的目标为$y_t=r_t+\gamma(1-\delta_t)\max_{a'}\overline Q(e_{\overline{\bm\theta}}(o_{t+1}),a',c)$，其中$\delta_t$是任务真正终止的标记。最小化$\tfrac12(Q(e_{\bm\theta}(o_t),a_t,c)-y_t)^2$时，梯度通过评论家进入编码器；目标网络在此次梯度中固定。重建损失可以辅助感知，但只有控制目标或充分性条件才能说明为什么该表示保留了动作选择所需信息。

在一个教学单步岔路中，目标左与目标右需要相反动作。若两项任务被编码为完全相同的输入，任何确定策略都只能满足其中一项；加上一个目标位即可消除这个表示冲突。另一方面，若左任务占采样的$90\%$，总体正确率$90\%$的常向左策略可能完全放弃右任务。这是任务权重造成的平均优势，不是两项任务都获得迁移。

训练计算随编码器前后向成本和样本数增长；多头设计还随任务头数量增加存储，模块化则增加组合与路由成本。诊断时应报告分任务回报、目标位消融与匹配容量的独立模型，区分信息被丢弃、容量不足和优化冲突。总体损失下降不能回答这些不同问题。

% Outline: 6.4.3
\subsection{可复用环境模型的训练}
\label{sec:rl-transfer-model-training}

共同地图中的不同风向提示一种分工：碰撞几何可以共享，偏移方向由任务控制。令$\widehat P_{\bm\theta,\bm\omega_c}(s'\mid s,a,\bm z_c)$预测下一状态，其中$\bm\theta$为共享参数，$\bm\omega_c$为可选的任务特定参数，$\bm z_c$为描述或历史推断出的任务变量。奖励头与观测头若也训练，应分别写作$\widehat R$和$\widehat O$；没有训练这些头就不能把模型称为包含全部反馈的模拟器。

以完全可观测的转移预测为例，取每条长度$H$的训练轨迹平均负对数似然：
\begin{equation}
 \mathcal L_{\text{模型}}
 =\mathbb E_{c\sim q}\mathbb E_{\tau\sim p_c^{\mu_c}}
 \left[-\frac1H\sum_{t=0}^{H-1}
 \log\widehat P_{\bm\theta,\bm\omega_c}
 (s_{t+1}\mid s_t,a_t,\bm z_c)\right].
 \label{eq:rl-transfer-model-loss}
\end{equation}
该目标来自最大化所采轨迹的条件转移似然；条件在已记录的$s_t,a_t$上，因而不把收集策略当成待拟合的动力学。不同长度回合按回合平均还是按步平均会赋予不同权重，必须选定一种。本式固定$H$，并以任务平均而非把全部记录直接混在一起定义目标。

图~\ref{fig:rl-transfer-model}~展示共同几何与任务条件模块如何组合。完全共享模型相当于忽略$\bm\omega_c,\bm z_c$；完全独立模型则为每项任务各拟合全部参数。模块化减少的是共同几何的重复估计，仍需估计新风向，不能以共享参数比例推算目标样本效率。

\begin{figure}[htbp]
  \centering
  % Original conditional model; final width 112 mm.
\begin{tikzpicture}[x=\linewidth,y=1mm,
 font=\figurefont\fontsize{8.5}{11}\selectfont,
 box/.style={draw=BookRule,fill=BookPaper,line width=.8pt,
   text width=24mm,minimum height=11mm,align=center,inner sep=1.5mm},
 flow/.style={-{Latex[length=2mm]},draw=BookInk,line width=.9pt}]
 \node[box] (sa) at (.14,16) {状态与动作\\$s_t,a_t$};
 \node[box] (shared) at (.51,16) {共享几何模块\\$\bm\theta$};
 \node[box] (context) at (.14,-6) {任务描述\\或历史};
 \node[box,fill=BookTeal!8] (task) at (.51,-6) {任务条件模块\\$\bm z_c,\bm\omega_c$};
 \node[box] (output) at (.87,5) {组合预测\\$\widehat P(s'\mid s,a,c)$};
 \draw[flow] (sa) -- (shared);
 \draw[flow] (context) -- node[above] {编码} (task);
 \draw[flow] (shared.east) -| (output.north);
 \draw[flow] (task.east) -| (output.south);
 \node[anchor=north,align=center,text=BookMuted] at (.5,-19)
   {完全共享：忽略任务条件\\完全独立：每项任务复制全部模块};
\end{tikzpicture}

  \caption{任务条件环境模型。状态动作进入共享模块，描述或历史进入任务模块，两者共同决定转移预测；奖励、观测头仅在确实训练时增加。图为本章模型参数化的原创示意，不表示所有环境都满足这种分解。}
  \label{fig:rl-transfer-model}
\end{figure}

一个可复核的教学例是前述一维风参数。在8次可辨认滑移的独立试验中观察到2次偏离，负对数似然为$-2\log\eta-6\log(1-\eta)$。令导数$-2/\eta+6/(1-\eta)=0$，得到$\widehat\eta=1/4$。随后把$\widehat\eta$送入共同碰撞模型，计算所有位置的转移；这一步依赖“同一个风参数适用于全图”的结构假设。

训练时知道$c$，部署时却只有历史，是一个额外推断问题。应训练从可见历史到$\bm z_c$的编码器，并在测试中只给这些历史；不能把真实风标签作为隐蔽输入。固定日志上的训练成本约为数据量乘单样本模型前后向成本，模型内规划还要乘候选动作序列数及前瞻深度。未覆盖风向、传感器变化和反复滚动的误差累积，都可能使训练似然较好的模型产生错误策略。

% Outline: 6.4.4
\subsection{可复用价值与策略的训练}
\label{sec:rl-transfer-value-training}

同一岔路在不同目标下需要不同动作，价值函数必须看见目标。目标条件MDP把目标$g$纳入任务描述，\term{通用价值函数逼近器}（universal value function approximator, UVFA）用一个$Q_{\bm\theta}(s,a,g)$同时概括状态与目标上的价值。它可以通过多目标TD目标训练：从任务$g$的行为策略采样，按该目标重算奖励和终止，拟合$r_t+\gamma(1-\delta_t)\max_{a'}\overline Q(s_{t+1},a',g)$。动作选择也必须查询相同的$g$，不能训练时条件化、执行时却丢掉它\scite{rl6-schaul2015-uvfa}

一个共享评论家可在网络中联合表示所有目标；多头输出保留共同特征、为每个已见任务独立输出价值；模块化策略则让不同模块处理几何、目标或控制模式。多头编号本身不解释新目标，UVFA的连续输入也不自动带来外推能力。单条TD更新至少要计算当前价值与目标动作值，有限动作枚举成本随$|\mathcal A|$增加，连续动作则需要行动者或其他最大化近似。

后继特征提供了更明确的价值复用条件。其前身\term{后继表示}（successor representation, SR）在固定策略下记录未来访问各状态的折扣次数。有限状态时可定义
\[
 M^\pi(s,u)=\mathbb E_\pi\!\left[\sum_{t\geq0}\gamma^t
       \mathbb I\{s_t=u\}\mid s_0=s\right].
\]
若$\bm P^\pi$为策略诱导的状态转移矩阵，则拆出$t=0$项得
$\bm M^\pi=\bm I+\gamma\bm P^\pi\bm M^\pi$，从而
$\bm M^\pi=(\bm I-\gamma\bm P^\pi)^{-1}$。逆矩阵存在是因为随机矩阵的谱半径不超过1且$\gamma<1$。该表示一般存储$|\mathcal S|^2$个数；将状态指示量换成低维事件特征，就得到SF\scite{rl6-dayan1993-successor}

\begin{definition}[固定策略的后继特征]
 \label{def:rl-transfer-sf}
 在共同动力学$P$、共同折扣$0\leq\gamma<1$与一致状态动作语义下，设
 $\sup_{s,a,s'}\lVert\bm\phi(s,a,s')\rVert_2\leq C_\phi<\infty$。
 固定平稳策略$\pi$的后继特征为
 \begin{equation}
 \bm\psi^\pi(s,a)=
 \mathbb E_\pi\!\left[
 \sum_{t=0}^{\infty}\gamma^t\bm\phi(s_t,a_t,s_{t+1})
 \,\middle|\,s_0=s,a_0=a\right].
 \label{eq:rl-transfer-sf}
 \end{equation}
 首动作按条件强制执行，之后服从$\pi$。
\end{definition}

现在完整连接这个对象与价值。对随机奖励，另假设各状态动作下的绝对一阶矩有共同有限上界；有界奖励样本是一个充分条件，从而折扣奖励和可积。由特征有界性，
$\sum_{t\geq0}\gamma^t\lVert\bm\phi_t\rVert_2\leq C_\phi/(1-\gamma)$；
对固定有限$\bm w_c$，条件期望奖励的绝对折扣和也有上界
$C_\phi\lVert\bm w_c\rVert_2/(1-\gamma)$。因此可以用支配收敛把部分和的极限移入期望，并用条件期望先把$r_t$替换成转移条件期望奖励。依次得到
\begin{align}
 Q_c^\pi(s,a)
 &=\lim_{N\to\infty}\mathbb E_\pi
       \left[\sum_{t=0}^{N}\gamma^t R_c(s_t,a_t,s_{t+1})
       \mid s_0=s,a_0=a\right]\notag\\
 &=\lim_{N\to\infty}\left\langle
       \mathbb E_\pi\left[\sum_{t=0}^{N}\gamma^t\bm\phi_t
       \mid s_0=s,a_0=a\right],\bm w_c\right\rangle\notag\\
 &=\langle\bm\psi^\pi(s,a),\bm w_c\rangle.
 \label{eq:rl-transfer-sf-value}
\end{align}
第二步使用式~\eqref{eq:rl-transfer-linear-reward}~与有限和、内积的线性性，第三步使用前述可积性。SF不含奖励权重，却仍依赖$P$和$\pi$；在新任务换掉转移后，原式右侧不再使用正确的轨迹分布。

拆出第一步特征，得到向量Bellman方程
\begin{equation}
 \bm\psi^\pi(s,a)=
 \mathbb E_{s'\sim P}\left[\bm\phi(s,a,s')
 +\gamma\mathbb E_{a'\sim\pi(\cdot\mid s')}\bm\psi^\pi(s',a')\right].
 \label{eq:rl-transfer-sf-bellman}
\end{equation}
因此可沿用价值评价的采样方法。表格情形从覆盖相关状态动作的行为策略$\mu$取得转移，计算
$\bm y_t=\bm\phi_t+\gamma\sum_{a'}\pi(a'\mid s_{t+1})\widehat{\bm\psi}(s_{t+1},a')$，
再令$\widehat{\bm\psi}(s_t,a_t)$向$\bm y_t$移动步长$\alpha$。终止时继续项为0。若仅采样一个$a'\sim\pi$，则以该样本替代动作期望。这里评价的是固定$\pi$，不能无说明地把下一动作改成最大值。表格每次更新的向量运算为$O(d|\mathcal A|)$，采样下一动作时为$O(d)$；网络实现还包含前后向成本。

新权重未知时，可从目标奖励样本最小化
$\sum_j(r_j-\bm\phi_j^\top\bm w_c)^2$；要唯一辨认权重，数据特征矩阵还需足够秩。若权重由目标描述直接给出，则不需要目标反馈拟合这一步，但仍需要在共同动力学下学好的SF。

已有多条策略时，可以在每个状态重新选动作，而不必整回合只运行某一条。\term{广义策略改进}（generalized policy improvement, GPI）先用同一目标奖励评价每条策略，再对它们的动作价值上包络取贪心。下面的保证比较已有策略，而不是未知的新任务最优策略\scite{rl6-barreto2017-successor}

\begin{theorem}[有一致估计误差的广义策略改进]
 \label{thm:rl-transfer-gpi}
 在同一个有限状态动作、有界奖励的折扣MDP中，给定有限平稳策略集
 $\{\pi_1,\ldots,\pi_m\}$。若所有状态动作均满足
 $|\widehat Q^{\pi_i}(s,a)-Q^{\pi_i}(s,a)|\leq\varepsilon$，
 定义
 \[
 \pi(s)\in\argmax_{a\in\mathcal A(s)}\max_i\widehat Q^{\pi_i}(s,a),
 \]
 并固定平局规则，则
 \begin{equation}
 Q^\pi(s,a)\geq\max_iQ^{\pi_i}(s,a)
                -\frac{2\varepsilon}{1-\gamma}.
 \label{eq:rl-transfer-gpi-bound}
 \end{equation}
\end{theorem}

证明思路的关键是把一次近似贪心损失变成可以递推的Bellman不等式。令
$f(s,a)=\max_iQ^{\pi_i}(s,a)$、$\widehat f(s,a)=\max_i\widehat Q^{\pi_i}(s,a)$；
一致误差保证$|f-\widehat f|\leq\varepsilon$。由$\pi$的贪心定义，
\[
 f(s,\pi(s))\geq\widehat f(s,\pi(s))-\varepsilon
 =\max_a\widehat f(s,a)-\varepsilon
 \geq\max_a f(s,a)-2\varepsilon.
\]
两次误差分别来自所选动作和最佳动作。用$T^\pi$表示动作价值的Bellman评价算子，便有
\begin{align*}
 T^\pi f
 &\geq R+\gamma P\max_a f-2\gamma\varepsilon\\
 &\geq \max_i\{R+\gamma P\mathbb E_{a'\sim\pi_i}Q^{\pi_i}\}
           -2\gamma\varepsilon\\
 &=f-2\gamma\varepsilon\geq f-2\varepsilon.
\end{align*}
其中$P$表示对后继状态取期望；第二行用“每个后继状态上的最大值至少覆盖任一$\pi_i$的动作平均”，再对策略取最大。Bellman单调性以及
$T^\pi(f-b)=T^\pi f-\gamma b$给出
$(T^\pi)^k f\geq f-2\varepsilon\sum_{j=0}^{k-1}\gamma^j$。
压缩性保证左侧收敛到$Q^\pi$，几何级数收敛到$1/(1-\gamma)$，于是得到式~\eqref{eq:rl-transfer-gpi-bound}。这一关系也说明深度网络上小的平均TD损失不等于定理要求的全状态动作一致误差。

\begin{example}[两条绕障策略的新奖励评价]
 \label{ex:rl-transfer-sf-routes}
 以下简化保留导航地图与固定吸收出口$(4,3)$，取消滑移与步成本，沿途地标允许重复计分。取特征为进入$(1,3)$、进入$(4,1)$、首次进入出口的三个指示量，吸收后特征为0，$\gamma=0.9$。两条确定策略从起点分别走“上、上、右、右、右”和“右、右、右、上、上”；离开这两条路线的状态可用任一固定合法规则补全策略。

 上路在$t=1$经过第一个地标，$t=4$进入出口；下路在$t=2$经过第二个地标，$t=4$进入出口。因此从起点执行各自首动作并继续原策略时，
 \[
 \bm\psi^{\pi_1}=(0.9,0,0.6561)^\top,\qquad
 \bm\psi^{\pi_2}=(0,0.81,0.6561)^\top.
 \]
 新权重$\bm w=(1,2,1)^\top$给出两条完整路线的回报$1.5561$与$2.2761$。固定调用下路优于固定调用上路；GPI则需为每个候选首动作查询每条补全策略的SF，再逐状态选择，不能仅凭这两个整路线数字实现。

 当SF精确时，GPI在各状态不会劣于库内任一策略；当所有动作价值一致误差至多为教学假设$\varepsilon=0.01$时，本定理给出的逐动作损失容限为$0.2$。从状态开始执行GPI也满足同样的保守容限：近似贪心至多损失$2\varepsilon$，后续至多损失$2\gamma\varepsilon/(1-\gamma)$。这些保证没有比较库外最优策略；地标重复奖励甚至可能使有益循环优于直接离场。
\end{example}

策略库有$m$条策略、每个SF为$d$维时，直接GPI查询约需$m|\mathcal A|$个SF输出和$O(m|\mathcal A|d)$次乘加。相比每个新奖励都从头学习，它复用了后果计算；相比单次策略前向，它增加了查询成本。新奖励残差、新动力学误差和策略库缺失各自需要额外分析，不能全塞进未经测量的$\varepsilon$。

另一种保存行为的方法是\term{策略蒸馏}（policy distillation）：在教师策略访问的状态上，让学生拟合教师动作分布。以共同动作接口为例，可最小化
\[
 \mathcal L_{\text{蒸馏}}=
 \mathbb E_{c\sim q,\,s\sim\nu_c}
 \operatorname{KL}\!\left(\pi_c^{\text{教师}}(\cdot\mid s)
       \,\Vert\,\pi_{\bm\theta}(\cdot\mid s,c)\right).
\]
$\nu_c$说明教师或其他行为收集的状态分布，梯度只更新学生。一个教学二动作教师给出$(0.8,0.2)$，学生给出$(0.5,0.5)$，则交叉熵对学生两个logit的梯度为$(-0.3,0.3)$，一次更新提高第一动作概率。压缩多个教师省去部署时多个网络的独立调用，但没有保存教师如何从失败中学习的过程，日志外状态仍可能失配。原始策略蒸馏工作以多任务策略压缩为主要用途\scite{rl6-rusu2016-distillation}

若部署允许少量梯度更新，训练目标可以直接考虑更新后的表现。\term{模型无关元学习}（model-agnostic meta-learning, MAML）学习一个适合后续更新的初始化。这里“模型无关”指适用于不同可微参数模型，不是强化学习中“无模型”的分类。令$L_c(\bm\theta)=-J_c(\pi_{\bm\theta})$为负回报，任务内更新算子为$U_c$，元目标写为
\begin{equation}
 \min_{\bm\theta}\ %
 \mathbb E_{c\sim p_{\text{训练}}}
 \mathbb E_{\mathcal D_c^{\text{内}}\sim p_c^{\pi_{\bm\theta}}}
 L_c\!\left(U_c(\bm\theta,\mathcal D_c^{\text{内}})\right).
 \label{eq:rl-transfer-meta-objective}
\end{equation}
外层抽任务，内层由当前策略采样并更新；更新后的策略还要在同一任务重新采集$\mathcal D_c^{\text{外}}$，估计更新后回报。外层不是让初始化立刻解决所有任务，而是让约定预算内的更新有用\scite{finn2017}

对固定内层数据和可微更新$\bm\theta'=\bm\theta-\alpha\nabla\widehat L_c^{\text{内}}$，链式法则给出条件导数
\[
 \nabla_{\bm\theta}L_c(\bm\theta')
 =\left(\bm I-\alpha\nabla_{\bm\theta}^2\widehat L_c^{\text{内}}\right)^\top
       \nabla_{\bm\theta'}L_c(\bm\theta').
\]
在强化学习中，内层数据分布本身依赖$\bm\theta$；若要求式~\eqref{eq:rl-transfer-meta-objective}~的完整梯度，还存在
$L_c(\bm\theta')\nabla_{\bm\theta}\log p_c^{\pi_{\bm\theta}}(\mathcal D_c^{\text{内}})$
这一采样分布项。把已采轨迹当作固定张量、对普通策略梯度代理反传，并不自动得到无偏的完整高阶梯度。必须说明使用的估计器及省略项；具体目标任务更新见第~\ref{sec:rl-transfer-gradient-adaptation}~节。

表~\ref{tab:rl-transfer-behavior-objects}~比较四种保存方式。技能还需保存启动集$I_o$、内部策略$\pi_o$及终止概率$\beta_o$；其发现机制见第~\ref{chap:rl-learning}~章。这里的关键区别是保存已有动作规律、可重加权的后果，还是未来更新的起点。

\begin{table*}[t]
 \centering\small
 \begin{tabularx}{\linewidth}{@{}>{\raggedright\arraybackslash}p{22mm}*{3}{>{\raggedright\arraybackslash}X}@{}}
  \toprule
  对象 & 保存什么 & 目标任务还需做什么 & 接口与边界\\
  \midrule
  价值与SF & 指定行为的长期后果 & 重加权、策略改进或价值重估 & 奖励类、共同转移、策略语义\\
  蒸馏策略 & 教师动作分布 & 条件调用或额外适应 & 输入与动作对齐；日志外误差\\
  可迁移初始化 & 适合更新的参数起点 & 用获准数据做有限更新 & 步长、梯度质量、更新后评价\\
  技能库 & 启动集$I_o$、输入、$\pi_o$、$\beta_o$ & 选择、组合或调整内部控制 & 动作效果、持续时间与终止含义兼容\\
  \bottomrule
 \end{tabularx}
 \caption{行为知识的四种保存方式。技能“能够运行”只验证部分接口，不保证适合新目标；表中没有把它们视为互斥算法类别。}
 \label{tab:rl-transfer-behavior-objects}
\end{table*}

% Outline: 6.5
\section{目标任务策略的形成}
\label{sec:rl-transfer-policy-formation}

共享模型或策略训练完毕，并不意味着目标任务策略已经确定。给定新目标可以直接查询条件策略；隐藏风向则需要从反馈中推断。新的信息可以存入记忆、任务后验或网络参数，是否允许继续收集信息又是另一项权限。本节先说明信息怎样形成，再分别讨论这些可以组合的机制。

% Outline: 6.5.1
\subsection{从显式描述与交互历史形成任务信息}
\label{sec:rl-transfer-belief}

考虑导航任务的一个独立教学变体：终点已知，风向$c\in\{\text{左风},\text{右风}\}$在整个试验期间固定，先验各为$1/2$。终点描述没有包含风向信息。在一个获准的测试位置，某个试探动作产生可见“右偏”的概率在两项任务中分别为$0.2$和$0.8$。观察一次右偏会改变对风向的判断；反复撞上同一堵墙，如果两种风都只留下“位置未变”，就可能完全没有区分力。数字仅用于说明似然更新，未作为基准网格的新转移解。

设当前物理状态可见，历史为
$h_t=(s_0,a_0,r_0,\delta_0,\ldots,s_t)$，其中$\delta_t$记录动作$a_t$之后的终止事件。有限任务集合上的\term{任务信念}（task belief）
$b_t(c)=\mathbb P(c\mid h_t)$是给定已见证据后对任务的概率分布。显式描述可以约束先验或直接确定某些任务分量；它与历史证据应进入同一条件化过程，而不是相互替代。若初始状态分布也随任务变化，看到$s_0$时便已有证据：描述条件下的先验为$p_0(c)$时，$b_0(c)\propto p_0(c)\rho_{0,c}(s_0)$，不能无条件沿用$p_0$。

令$\ell_c(s_{t+1},r_t,\delta_t\mid s_t,a_t)$为任务$c$中可用模型给出的联合似然。任务固定、奖励及终止模型正确，且动作仅由已记录历史和独立随机数选择时，贝叶斯公式给出
\begin{align}
 b_{t+1}(c)
 &=\frac{b_t(c)\,
     \ell_c(s_{t+1},r_t,\delta_t\mid s_t,a_t)}
 {\sum_{\widetilde c\in\mathcal C}b_t(\widetilde c)\,
     \ell_{\widetilde c}(s_{t+1},r_t,\delta_t\mid s_t,a_t)}.
 \label{eq:rl-transfer-belief-update}
\end{align}
分母为所见反馈的预测概率，需大于0。推导中，联合概率包含
$b_t(c)\pi(a_t\mid h_t)\ell_c(\cdot)$；动作概率在各候选任务间相同，故在归一化时消去。若收集者另知隐藏任务并据此选动作，却未记录这项信息，就不能不加分析地消去它。

前述右偏观测给出
$b_{t+1}(\text{右风})=0.8/(0.8+0.2)=0.8$。若独立重置到测试位置后又观察一次右偏，则后验为
$0.8^2/(0.8^2+0.2^2)=16/17$。任务后验改变的是对固定对象的认识，不是风向本身。若物理状态也隐藏，应维护
$b_t(s,c)=\mathbb P(s_t=s,c\mid h_t)$，先对旧状态求和预测，再乘观测与奖励似然；不能用一个任务向量代替状态滤波而省去所需条件。

精确枚举有限任务的单步更新需要$|\mathcal C|$次似然计算；状态同时隐藏时还包含状态转移求和。大任务空间通常改用历史编码器近似推断。图~\ref{fig:rl-transfer-belief-loop}~将推断和控制组成闭环，虚线单独标出训练时的监督或变分目标，避免把网络输出误画成真实任务标签。

\begin{figure}[htbp]
 \centering
 % Original information loop; final width 112 mm. No empirical data.
\begin{tikzpicture}[x=\linewidth,y=1mm,
 font=\figurefont\fontsize{8.5}{11}\selectfont,
 box/.style={draw=BookRule,fill=BookPaper,line width=.8pt,
   text width=28mm,minimum height=12mm,align=center,inner sep=1.5mm},
 flow/.style={-{Latex[length=2mm]},draw=BookInk,line width=.9pt},
 lab/.style={fill=white,inner sep=1mm}]
 \node[box] (history) at (.19,13) {可见历史\\$h_t$};
 \node[box,fill=BookTeal!8] (belief) at (.80,13)
   {任务信息\\$b_t$或$q(\bm z\mid h_t)$};
 \node[box] (experience) at (.19,-17) {新经验\\观测、奖励、终止};
 \node[box] (policy) at (.80,-17) {条件策略\\选择动作};
 \node[box,text width=42mm,fill=BookAmber!8] (train) at (.49,43)
   {仅训练：预测、价值信号\\或变分目标};
 \draw[flow] (history) -- node[above] {编码或更新} (belief);
 \draw[flow] (belief) -- node[right] {条件输入} (policy);
 \draw[flow] (policy) -- node[below] {环境交互} (experience);
 \draw[flow] (experience) -- node[left] {追加} (history);
 \draw[flow,dashed,draw=BookAmber] (train.east) -| (belief.north);
 \node[anchor=north,align=center,text=BookMuted] at (.50,-29)
   {潜变量不是自动已知的任务标签\\执行时仅用获准描述与反馈};
\end{tikzpicture}

 \caption{任务信息与行为的闭环。实线为执行时可见的信息流，虚线为训练时用于塑造推断器的目标；“任务信息”可以是精确后验、近似分布或普通记忆，三者不能仅因向量维数相同就互换。图为本章机制的原创归纳。}
 \label{fig:rl-transfer-belief-loop}
\end{figure}

PEARL用转移上下文$\mathcal C_t$形成近似分布$q_{\bm\varphi}(\bm z\mid\mathcal C_t)$，再从中抽样$\bm z$，交给条件策略与评论家。其原始设计把各条上下文对应的高斯因子组合起来，适合任务固定、转移可按集合处理的情形；离策略回放用于训练评论家与策略，上下文编码器由价值学习信号及相对先验的KL正则约束。训练抽任务、抽该任务上下文、抽回放转移，更新相应损失；部署则收集允许的少量转移、更新潜变量分布并再抽样行为。这里保存的是近似任务信息，不是已知真实风标签\scite{rl6-rakelly2019-pearl}

VariBAD也近似任务后验，但通过循环编码历史与变分目标学习潜变量：以潜变量重建同一任务的奖励和转移，并用KL项约束后验与先验或先前后验。条件策略读取后验的参数，而不只读取一个任务点估计，训练时以跨回合回报学习兼顾信息获取与利用。其执行阶段不需要调用训练用解码器做模型内规划。因此，“有预测解码器”不等于“部署时基于模型规划”\scite{rl6-zintgraf2020-varibad}

两种设计的取舍可以在同一个右偏例子中检查。精确二任务模型给出$0.8$这一可核后验；近似编码器则需检验同类历史的预测与校准，不能要求某个潜变量坐标也恰好等于$0.8$。PEARL式抽样让一次行为在某个任务假设下连贯执行，VariBAD式条件策略可利用分布的不确定性决定是否再探测；它们是否学到有效探索仍依赖训练任务与反馈。前者要计算上下文因子及条件策略，后者要逐步维护循环状态和后验参数；两者训练还包含评论家或解码器成本，不能只比较策略前向次数。

存在比近似误差更根本的限制。构造两个终点隐藏的教学任务：在最后一次岔路选择前，预算内所有获准动作的状态、奖励和终止分布完全相同；最终向左或向右分别只在其中一个任务成功，先验各为$1/2$。任何历史编码器收到的数据分布相同，因此最终选择左的概率在两任务中相同。若该概率为$p$，平均成功率为$\tfrac12p+\tfrac12(1-p)=1/2$。增加网络容量不能提高这项上限。只有允许额外访问会显示终点线索的状态，或提供带区分力的描述，主动辨识才可能改变信息条件。

% Outline: 6.5.2
\subsection{无目标任务更新数据下的零样本策略}
\label{sec:rl-transfer-zero-shot}

目标坐标已给出、输入语义与训练一致时，策略形成可以只是一次条件查询：
$a_t\sim\pi_{\bm\theta}(\cdot\mid s_t,g)$，或
$a_t\in\argmax_a Q_{\bm\theta}(s_t,a,g)$。若同时给出可解释的新奖励权重，还可计算
\[
 a_t\in\argmax_a\max_i
 \langle\widehat{\bm\psi}^{\pi_i}(s_t,a),\bm w_g\rangle.
\]
后一种调用复用了SF策略库，不需要目标奖励样本；但共同动力学、终止与特征语义，以及第~\ref{thm:rl-transfer-gpi}~定理的估计条件仍不能省略。

这里沿用第~\ref{sec:rl-transfer-budget}~节的零样本口径：禁止用目标反馈更新权重、任务记忆或选择检查点，无论更新发生在回合内还是跨回合。当前观测仍然可见，否则连位置控制都无法执行；预先约定的物理状态滤波与未知任务辨识应分开报告。若把偏移或奖励反馈送入历史编码器，据此推断风向并改变后续控制，就已发生任务适应，即使尚未结束第一个回合、也没有一次反向传播。

可以在同一训练模型上固定目标坐标格式、动作选择规则和推理次数，分别测试图~\ref{fig:rl-transfer-task-plane}~中的$(1,0.15)$、$(3,0.3)$和$(4,0.45)$任务。它们依次代表支持内新任务、未见因素组合和范围外因素；训练参数与测试权限完全相同，变化的是所检验的任务关系。应分别报告三组零样本回报，而不是把它们合成一个“新任务平均分”。

若输入只有新编号3，条件网络无法由数字推断它代表哪个终点。若新目标需要从未形成的控制方式，或目标设备交换了左右动作接口，旧策略库即使每条策略都能前向计算，也可能没有一个合法有效的候选。零样本能力来自已有知识与描述之间成立的关系，不来自关闭更新开关本身。

% Outline: 6.5.3
\subsection{参数固定条件下基于上下文的策略适应}
\label{sec:rl-transfer-context-adaptation}

同一隐藏风向任务连续运行两个回合，第一回合先试探，第二回合依据偏移结果绕障，这已经是适应。变化可以只发生在记忆中。将“学习”等同于“修改网络权重”，会漏掉这种由历史改变决策规则的机制。

RL$^2$以循环策略把快速学习过程放入内部状态\scite{rl6-duan2016-rl2}
执行时固定网络参数$\bm\theta$，令
\begin{align}
 \bm m_t
 &=f_{\bm\theta}(\bm m_{t-1},o_t,a_{t-1},r_{t-1},\delta_{t-1}),\notag\\
 a_t&\sim\pi_{\bm\theta}(\cdot\mid\bm m_t).
 \label{eq:rl-transfer-rl2}
\end{align}
首步的上一动作、奖励与终止标记使用约定的空值。终止标记说明上一动作结束了物理回合；环境此时重置位置，但同一任务的$\bm m$继续保留。只有切换任务、开始新的独立适应试验时，才重置任务记忆。训练的慢过程在多项任务上优化跨多个回合的总回报，使早期收集信息可能由后续收益补偿。

图~\ref{fig:rl-transfer-rl2-memory}~把固定参数外框与可变记忆分开。在前述教学风任务中，可构造一个实现贝叶斯更新的记忆器：初始记忆存右风概率$1/2$，第一个回合观察右偏后存$0.8$，第二回合用新概率选择路线。这个例子证明固定参数可以实现适应，不声称训练出的RL$^2$必然等价于精确贝叶斯推断。

\begin{figure}[htbp]
 \centering
 % RL^2 protocol redrawn from Duan et al. (2016), arXiv:1611.02779v2.
% Final width 112 mm; theta is fixed at deployment, memory is not.
\begin{tikzpicture}[x=\linewidth,y=1mm,
 font=\figurefont\fontsize{8.5}{11}\selectfont,
 box/.style={draw=BookRule,fill=white,line width=.8pt,
   minimum height=12mm,align=center,inner sep=1.5mm},
 flow/.style={-{Latex[length=2mm]},draw=BookInk,line width=.9pt}]
 \draw[draw=BookTeal,line width=.8pt,fill=BookTeal!5]
   (.28,-24) rectangle (.98,23);
 \node[anchor=north] at (.63,21) {部署参数$\bm\theta$固定};
 \node[box,text width=23mm] (input) at (.115,0)
   {当前$o_t$\\前一步$a_{t-1}$\\$r_{t-1},\delta_{t-1}$};
 \node[box,text width=23mm] (memory) at (.49,0)
   {可变循环记忆\\$\bm m_t$};
 \node[box,text width=17mm] (action) at (.84,0)
   {策略输出\\$a_t$};
 \draw[flow] (input) -- (memory);
 \draw[flow] (memory) -- (action);
 \draw[flow] (memory.south) -- (.49,-16) -- (.34,-16)
   |- (memory.west);
 \node[anchor=west] at (.53,-16) {供下一步使用};
 \node[anchor=north,align=center] at (.50,-30)
   {回合结束：重置环境，保留同任务记忆\\
    任务切换：重置记忆，开始独立适应试验};
\end{tikzpicture}

 \caption{固定参数中的可变记忆。输入包含当前观测及上一动作对应的奖励和终止标记；同任务回合结束仅重置物理环境，任务改变才重置记忆。图据RL$^2$的交互协议重绘，贝叶斯数值例是本章构造，不是其训练结果。}
 \label{fig:rl-transfer-rl2-memory}
\end{figure}

循环记忆并不是唯一实现。Algorithm Distillation（AD，算法蒸馏）从源强化学习算法的多回合学习记录中抽取长序列，以自回归动作预测拟合
$-\sum_t\log\pi_{\bm\theta}(a_t\mid h_t)$\scite{rl6-laskin2023-ad}
输入历史包括观测、动作与奖励，训练样本保留从探索到改进的顺序；部署时把自身新经历追加到历史，固定参数继续预测动作。它模仿的是学习过程中的行为变化，而前述策略蒸馏通常只压缩已训练教师的行为。

Decision Pretrained Transformer（DPT，决策预训练Transformer）则以任务内上下文和查询状态预测教师最优动作。训练时从任务$c$抽取上下文$\mathcal D_c$和查询$s$，拟合
$-\log\pi_{\bm\theta}(a_c^*(s)\mid s,\mathcal D_c)$；连续动作可使用相应密度损失。理想化地，若标签确为各任务最优动作、任务分布与上下文生成方式正确覆盖测试情形，监督条件分布会对仍与上下文相容的任务作后验混合。这解释了与后验采样的联系，不是任意有限网络都具备贝叶斯最优探索的保证\scite{rl6-lee2023-dpt}

\Needspace{26\baselineskip}
另一些方法直接以强化学习训练长期记忆。Adaptive Agent（AdA）在程序化任务空间与课程中训练带记忆的智能体，让跨回合反馈影响行为；其训练信号是交互回报，而非给定最优动作标签\scite{rl6-bauer2023-ada}

Scalable In-Context Q-Learning（S-ICQL）把上下文条件化放进价值学习：利用世界模型生成任务提示，训练上下文条件的$Q/V$，结合偏向高价值样本的expectile目标和优势加权策略提取；测试上下文可以来自规定的在线交互或离线记录。

Vintix II沿用DPT式决策预训练思路，以多领域、含噪行为记录构造上下文，并用示范者动作重标记监督目标；连续动作生成采用流匹配目标。它同时研究在线与离线的上下文使用。此处的跨领域结果仍受所用低维观测、控制接口和训练任务支持约束，不能改写成任意视觉任务的通用保证。

表~\ref{tab:rl-transfer-context-methods}~按信息与策略形成方式对照这些方法。它们都能在固定参数下利用新上下文，但训练信号分别来自交互回报、学习历史中的动作、教师标签或条件价值估计。这些差别决定了模型能够从训练数据中学到什么，不能仅以是否使用Transformer判断适应机制。

\begin{table*}[t]
 \centering\small
 \begin{tabularx}{\linewidth}{@{}>{\raggedright\arraybackslash}p{20mm}*{4}{>{\raggedright\arraybackslash}X}@{}}
  \toprule
  方法 & 上下文内容 & 训练数据来源 & 目标反馈 & 策略提取\\
  \midrule
  RL$^2$ & 观测、前动作、奖励、终止 & 多任务在线跨回合交互 & 同任务内逐步到达 & 循环记忆条件策略\\
  AD & 按顺序的多回合学习历史 & 源RL算法的学习轨迹 & 追加自身交互历史 & 自回归预测动作\\
  DPT & 任务内转移及查询状态 & 上下文与教师最优动作标签 & 给定上下文或新交互 & 监督条件动作分布\\
  AdA & 跨回合观测、动作与反馈 & 程序化任务、课程与在线RL & 试验内持续反馈 & 记忆条件的RL策略\\
  S-ICQL & 任务提示与交互上下文 & 世界模型提示和离线记录 & 在线或离线上下文 & 条件价值与优势加权策略\\
  Vintix II & 多领域转移上下文 & 含噪记录、示范者动作重标 & 在线或离线上下文 & DPT式条件化与流匹配动作头\\
  \bottomrule
 \end{tabularx}
 \caption{固定参数上下文适应的不同训练信号。各方法的上下文、反馈量和策略提取成本必须单独核对；该表归纳原始方法，不表示使用相同数据得到的实验排名。}
 \label{tab:rl-transfer-context-methods}
\end{table*}

表中的策略提取方式也影响上下文如何用于决策。AD拟合学习历史中的动作，S-ICQL则经由条件价值评价动作；其中expectile目标对正负残差赋予不同的二次损失权重，不是直接把最大样本当作真实最优值。二者都可在固定参数下适应，却不能仅按“是否预测下一个动作”视为同一种学习机制\scite{rl1-liu2026-sicql}

这类方法的运行代价不止权重存储。循环策略每步更新固定维度记忆，长序列Transformer在朴素实现中对长度$L$的完整注意力计算为$O(L^2)$；缓存可以减少重复计算，但仍要存储上下文及相关状态。条件价值可能需要额外候选动作评价；Vintix II采用的流匹配动作头还需若干生成步骤\scite{rl1-polubarov2026-vintix}
比较时应报告实际上下文长度、缓存、推理次数与延迟，不把“零梯度步”写成“零适应成本”。

任务边界错误也会改变结果：把前一任务记忆带入下一任务会污染推断，过早清空记忆会丢掉跨回合证据；截断窗口可能删掉唯一的风向线索。学习轨迹若只有无效随机行为，AD无法仅靠动作模仿获得未出现的改进规则；教师标签、世界模型和任务生成器也各有支持限制。训练中学到某种改进模式，并不意味着在任意新任务上都能保证探索和性能单调提升。

% Outline: 6.5.4
\subsection{允许有限梯度更新下的参数适应}
\label{sec:rl-transfer-gradient-adaptation}

若协议允许用目标任务数据改动参数，复用初始化就成为一次有明确数据依赖的学习过程。设任务$c$在适应期间固定，从$\bm\theta^{(0)}=\bm\theta$出发，以当前策略收集$\mathcal D_c^{(0)}$。对于有限时域折扣回报，令
$L_c(\bm\theta)=-\mathbb E[\sum_{t=0}^{H-1}\gamma^t r_t]$；
终端价值取0。沿用策略梯度的可微性、环境不依赖策略参数及可积性条件，可构造
\begin{align}
 \widehat{\bm g}_c
 &=-\frac1N\sum_{j=1}^N\sum_{t=0}^{H-1}
   \gamma^t\nabla_{\bm\theta}\log\pi_{\bm\theta}
        (a_{j,t}\mid h_{j,t})
   \bigl(G_{j,t}-b(h_{j,t})\bigr),\notag\\
 \bm\theta'&=\bm\theta-\alpha\widehat{\bm g}_c.
 \label{eq:rl-transfer-inner-update}
\end{align}
$G_{j,t}$是轨迹$j$从动作时刻$t$起的折扣回报，基线$b$不能依赖当前所采动作，且在此梯度中视为固定。式中额外的$\gamma^t$来自起点目标的折扣；损失是负回报，所以减去负回报梯度才对应提高回报。该样本梯度可记作
$\nabla_{\bm\theta}\widehat L_c$，但不能通过把环境奖励当成参数的普通可微函数来获得它。

\begin{example}[一次目标任务参数更新]
 \label{ex:rl-transfer-gradient-step}
 将导航末端岔路简化为独立单步教学任务：向右奖励1、向左奖励0，动作概率为
 $p_\theta=\pi_\theta(\text{右})=\sigma(\theta)$，初值$\theta=0$，基线0。收集4次独立试验，假设实际抽到3次右、1次左。对于右动作，
 $\partial_\theta\log p_\theta=1-p_\theta=1/2$；左动作奖励0，贡献为0。于是
 \[
 \begin{aligned}
 \widehat g&=-\frac{3(1/2)}4=-\frac38,\qquad
 \alpha=2,\qquad \theta'=\frac34,\\
 p_{\theta'}&=\frac1{1+e^{-3/4}}\approx0.6792.
 \end{aligned}
 \]
 更新前的真实期望回报为$0.5$，更新后为$0.6792$。这里能够直接算真实值，是因为教学任务的奖励模型已完整给出；实际未知任务需要另采评价数据。4条记录的经验回报$0.75$既不是更新前真实值，也不能作为更新后性能证明。
\end{example}

执行顺序因此是收集、估计、更新、另行评价。扩展为$K$步时，用
$\bm\theta^{(k+1)}=\bm\theta^{(k)}-\alpha_k\widehat{\bm g}_c^{(k)}$；
每步若按同策略估计，需要新策略的数据，复用旧数据则必须采用适当离策略估计并满足覆盖条件。固定日志不提供未记录动作的反事实回报，允许梯度计算也不会自动补齐这一缺口。

第~\ref{sec:rl-transfer-value-training}~节的MAML外循环还需评价更新对初始化的依赖。保留
$\bm I-\alpha\nabla^2\widehat L_c$是在给定内层样本时反传内循环导数；一阶近似以$\bm I$代替这一Jacobian，减少高阶自动微分开销。二者都还需交代强化学习采样分布项是否纳入估计，不能只因使用了二阶自动微分就声称得到精确随机元梯度。部署时执行若干普通目标梯度步，本身并不需要再计算元梯度。

表~\ref{tab:rl-transfer-storage}~按新信息写入的位置比较适应机制。计算成本包括训练所需的反向传播、优化器状态和独立评价，而不只是最终策略推理。

\begin{table*}[t]
 \centering\small
 \begin{tabularx}{\linewidth}{@{}>{\raggedright\arraybackslash}p{19mm}*{4}{>{\raggedright\arraybackslash}X}@{}}
  \toprule
  机制 & 信息存储位置 & 变化内容 & 重置与可逆性 & 主要代价\\
  \midrule
  上下文适应 & 历史窗口、循环状态 & 新经验与压缩记忆 & 可清空；截断丢失的信息未必可恢复 & 记忆存储、序列前向\\
  任务信念 & 离散概率或后验参数 & 各假设的相对可信度 & 可重置先验；恢复旧后验需保存证据或快照 & 似然、归一化或近似推断\\
  参数适应 & 策略、价值或模型权重 & 梯度及优化器状态 & 有初值快照才可完整回滚 & 前后向、优化器、再评价\\
  \bottomrule
 \end{tabularx}
 \caption{新任务信息的存储位置。这是机制比较而不是互斥分类：后验可作为上下文输入，记忆更新也可与参数更新同时执行。}
 \label{tab:rl-transfer-storage}
\end{table*}

少样本梯度方差大，过大步长会破坏原有行为；模型错配或离策略权重过大还会使估计方向失真。教学例中的一次改进不能推成一般单调保证。没有目标更新数据、只获准固定参数部署，或探测动作违反可行动作约束时，就不能默认运行这套流程。

% Outline: 6.5.5
\subsection{持续交互下的任务辨识与在线适应}
\label{sec:rl-transfer-identification}

未知风向固定不变时，早先的有效观测仍然指向当前任务。持续交互可以逐渐减少不确定性，但前提是所选行为确实带来区分性证据。一直沿安全边界走、从不经过受风影响的位置，可能取得稳定回报，却永远无法辨认风向。

把这一取舍简化为独立的两阶段、不折扣教学任务。两种风先验各半；终端安全路线奖励$0.6$，按正确风向选择的捷径奖励1、选错奖励0。直接选择时，捷径期望$0.5$，故安全路线较好。若第一阶段允许付出$0.1$奖励成本，执行一个正确率$0.8$的对称风向探测，再按信号选择捷径，则总期望回报为$-0.1+0.8=0.7$。不探测的策略在第一阶段等待且不扣分，两者使用相同两阶段时域。探测提高了终端决策的期望回报$0.2$，付出成本$0.1$，净收益$0.1$；若探测成本改为$0.3$，原方案就不再值得。

这说明信息收益必须通过后续决策评价，不能仅以任务分类准确率评价。在线过程可以是每收集一小批转移便更新一次后验和策略，也可以逐步更新记忆、偶尔更新参数。共同结构是当前任务信息决定行为，行为决定新数据，新数据再更新任务估计或控制规则。只更新后验、只更新参数、同时更新两者，都可能形成闭环，但代价和稳定性不同。

图~\ref{fig:rl-transfer-identification-curves}~刻意把任务辨识误差和控制回报画在两个坐标系中。早期误差较大，策略改进受到信息不足限制；后来误差已经很低，回报仍未达到目标，提示还需改善规划或控制。图中的解析曲线仅用于区分两种诊断，不是任何方法的收敛率，也没有假设分类误差可以直接换算成回报。

\begin{figure}[htbp]
 \centering
 % Teaching analytic functions, not measured data or a convergence guarantee.
% B in [0,80]; e(B)=.5 exp(-B/8); J(B)=.2+.8(1-exp(-max(B-20,0)/16)).
% B=20 separates explanatory regions only. Final width 112 mm.
\begin{tikzpicture}
 \begin{axis}[name=ident,width=.98\linewidth,height=44mm,
   font=\figurefont\fontsize{8.5}{11}\selectfont,
   axis lines=left,xmin=0,xmax=80,ymin=0,ymax=.55,
   xtick={0,20,40,60,80},xticklabels={},
   ytick={0,.25,.5},ylabel={辨识误差$e(B)$},
   tick label style={font=\figurefont\fontsize{8}{10}\selectfont}]
  \path[fill=BookBlue!7] (axis cs:0,0) rectangle (axis cs:20,.55);
  \draw[BookRule,dashed] (axis cs:20,0) -- (axis cs:20,.55);
  \addplot[BookBlue,line width=1pt,domain=0:80,samples=81] {.5*exp(-x/8)};
  \node[align=center] at (axis cs:12,.40) {信息\\不足};
  \node[align=center] at (axis cs:49,.30) {后期辨识误差已低};
 \end{axis}
 \begin{axis}[at={(ident.south west)},anchor=north west,yshift=-7mm,
   width=.98\linewidth,height=44mm,
   font=\figurefont\fontsize{8.5}{11}\selectfont,
   axis lines=left,xmin=0,xmax=80,ymin=0,ymax=1.05,
   xtick={0,20,40,60,80},ytick={0,.5,1},
   xlabel={累计目标交互$B$（步）},ylabel={策略回报$J(B)$},
   tick label style={font=\figurefont\fontsize{8}{10}\selectfont}]
  \path[fill=BookBlue!7] (axis cs:0,0) rectangle (axis cs:20,1.05);
  \draw[BookRule,dashed] (axis cs:20,0) -- (axis cs:20,1.05);
  \addplot[BookTeal,line width=1pt,domain=0:80,samples=81]
    {.2+.8*(1-exp(-max(x-20,0)/16))};
  \node[align=center] at (axis cs:12,.67) {控制受\\信息限制};
  \node[align=center] at (axis cs:49,.38) {已辨识\\控制尚未优化完毕};
 \end{axis}
\end{tikzpicture}

 \caption{固定任务下辨识与控制的不同进度。教学解析曲线为$e(B)=0.5e^{-B/8}$和$J(B)=0.2+0.8(1-e^{-\max(B-20,0)/16})$；$B$为累计目标交互步数。竖线$B=20$用于分区说明，非实测拐点，亦非理论阈值。两图共享横轴，不共享纵轴量纲。}
 \label{fig:rl-transfer-identification-curves}
\end{figure}

误差可以按已知测试标签计算，也可以按后验对真实任务的概率或参数偏差计算，报告时需固定一种定义。部署系统没有真实标签时，只能检查预测残差、校准或其他可观察证据；不能把测试者知道的风向反送给策略。实际测量还应改变探测预算，分别检查信息是否改善，以及相同信息下控制优化是否改善。

探测、回到测试位置、失败恢复和再评价全部消耗资源。更多交互只有在可识别、能到达相关状态且反馈权限充分时才有帮助。若辨识需要越过禁止区域，不能以“学习所必需”为理由默认开放动作；安全探索与约束可行性见第~\ref{chap:rl-trustworthy}~章。

% Outline: 6.5.6
\subsection{任务持续变化下的跟踪与适应}
\label{sec:rl-transfer-tracking}

现在把固定风向改成回合间突变，或让风强度缓慢漂移。上个回合的右偏仍是真实记录，却可能不再代表当前风向。直接把式~\eqref{eq:rl-transfer-belief-update}~用于全部历史，会把已失效的证据累积成过度自信；要继续作贝叶斯推断，需要额外给出任务变化模型，否则应采用明确的跟踪估计。

变化检测与重启先检查预测残差是否异常，再丢弃或下调旧模型；滑动窗口只保留最近$W$条记录；历史降权给年龄为$j$的样本权重$\lambda^j$，其中$0<\lambda<1$。窗口大或$\lambda$接近1时，稳定阶段估计方差较低，变化后的响应却较慢。例如独立Bernoulli观测在固定参数下，窗口均值的方差为$\eta(1-\eta)/W$；参数突变后，这个均值在旧记录移出前仍混合两种机制。检测阈值也要在误报、漏报和延迟间取舍。

为明确比较对象，考虑共同有限状态动作空间、共同初始分布$\rho_0$的回合式任务序列。第$k$回合的MDP为$\mathcal M_k$，在该回合内固定，时域为$H$；本节用不折扣总奖励
$V_{k,0}^{\pi}(\rho_0)=\mathbb E_{\rho_0,\pi,\mathcal M_k}
[\sum_{t=0}^{H-1}r_{k,t}]$，终端价值0。$r_{k,t}$仍按回合内动作时刻编号，与无限折扣目标分别使用。令$\pi_k$是算法基于此前信息选择的本回合行为规则，可以在回合内依靠新历史行动。

\begin{definition}[动态遗憾与静态遗憾]
 \label{def:rl-transfer-regret}
 对$T$个回合，动态遗憾取每回合知道当前MDP、但不知道未来随机结果的最优策略作比较：
 \begin{equation}
 \operatorname{DReg}_T
 =\sum_{k=0}^{T-1}\left[
 V_{k,0}^{*}(\rho_0)-V_{k,0}^{\pi_k}(\rho_0)\right].
 \label{eq:rl-transfer-dynamic-regret}
 \end{equation}
 相对于预先规定的固定策略类$\Pi_{\text{固定}}$，静态遗憾为
 \begin{equation}
 \operatorname{SReg}_T
 =\max_{\pi\in\Pi_{\text{固定}}}
       \sum_{k=0}^{T-1}V_{k,0}^{\pi}(\rho_0)
       -\sum_{k=0}^{T-1}V_{k,0}^{\pi_k}(\rho_0).
 \label{eq:rl-transfer-static-regret}
 \end{equation}
 固定比较策略不随回合改变，可依赖回合内时间与允许观测；不能暗含当回合隐藏任务标签。随机任务序列或随机算法下，另对这些随机性取期望。
\end{definition}

动态比较者在每次变化后都立即知道新的奖励与转移，显著强于部署者。静态比较者只选一个共同规则；算法若成功跟踪变化，甚至可能优于它，故静态遗憾不应解释为适应延迟。两者必须使用相同初始分布；如果还改变初始分布，应把它列入变化机制和评价条件。

次线性动态遗憾需要约束环境变化或赋予额外信息。例如可限制相邻回合任务改变的次数
$S_T=\sum_{k=1}^{T-1}\mathbb I\{\mathcal M_k\ne\mathcal M_{k-1}\}$，
或者在共同有界奖励尺度下定义
\begin{align}
 B_R&=\sum_{k=1}^{T-1}\sup_{s,a}|R_k(s,a)-R_{k-1}(s,a)|,\notag\\
 B_P&=\sum_{k=1}^{T-1}\sup_{s,a}
 \lVert P_k(\cdot\mid s,a)-P_{k-1}(\cdot\mid s,a)\rVert_1.
 \label{eq:rl-transfer-variation}
\end{align}
变化次数、累计幅度与单次最小可检测幅度不是同一个量。频繁微小漂移可以次数多而累计幅度小；一次大跳变也可能因相关状态不可达而难以检测。是否给出变化提示、提示是否含新任务参数，以及是否已知变化预算，都需单独说明。相应非平稳学习理论是在这些条件下研究性能，而非取消所有条件后承诺及时适应\scite{rl6-wei2021-nonstationary}

\begin{theorem}[无预告快速变化的两臂下界]
 \label{thm:rl-transfer-linear-regret}
 在$H=1$的教学任务族中，每轮$k$独立均匀抽取
 $C_k\in\{1,2\}$，动作集合为两个臂，奖励为
 $r_k(a)=\mathbb I\{a=C_k\}$。选择动作前不给任何关于$C_k$的线索；选择后可以观察所得奖励。对任何使用过去历史并可随机化的算法，
 \[
 \mathbb E[\operatorname{DReg}_T]=T/2.
 \]
\end{theorem}

\begin{proof}
 令$\mathcal F_k$包含第$k$轮之前的全部动作、奖励及算法已用的随机数。算法还可使用与任务序列独立的新随机数$U_k$，据
 $(\mathcal F_k,U_k)$选择$A_k$。由构造，$C_k$与这两项信息独立且均匀，因此不论条件动作概率如何，
 对条件概率为正的动作$a$，
 $\mathbb P(C_k=a\mid\mathcal F_k,U_k,A_k=a)=1/2$。
 对动作分支求和便得
 \begin{align*}
 \mathbb P(A_k=C_k\mid\mathcal F_k,U_k)
 &=\sum_{a=1}^2
    \mathbb P(A_k=a\mid\mathcal F_k,U_k)\,\tfrac12\\
 &=\tfrac12.
 \end{align*}
 取全期望，每轮算法的期望奖励为$1/2$。知道当前任务$C_k$的动态比较者选择臂$C_k$，每轮奖励恒为1；它不需要预知任何未来任务。由期望的线性性，
 \[
 \mathbb E[\operatorname{DReg}_T]
 =\sum_{k=0}^{T-1}\left(1-\mathbb E[r_k(A_k)]\right)
 =\sum_{k=0}^{T-1}(1-\tfrac12)=T/2.
 \]
 轮末反馈即使揭示了本轮完整任务，也不能预测独立的下一轮，所以结论不依赖算法的记忆容量、计算时间或参数更新能力。
\end{proof}

这一下界不否认缓慢变化可以跟踪，也不否认预告信号可以帮助决策。它只指出：比较者知道当前最好动作而算法在选择前毫无线索时，无规律快速变化本身足以产生线性代价。静态比较者的期望优势却小得多。令$N_1$为$T$轮中臂1高奖励的次数，最佳固定臂奖励为
$T/2+|N_1-T/2|$；由方差及Cauchy--Schwarz不等式，
$\mathbb E|N_1-T/2|\leq\sqrt{T}/2$。因此同一算法对固定臂的期望静态遗憾至多$\sqrt{T}/2$，而动态遗憾仍是$T/2$。换比较对象就换了结论。

图~\ref{fig:rl-transfer-tracking}~进一步区分变化发生、变化检测和性能恢复三个时刻。检测准确也不等于立即恢复：模型重估、规划与控制学习仍有代价。恢复阈值和观察窗口应提前规定，不能看到曲线后再选最有利时刻。

\begin{figure}[htbp]
 \centering
 % Teaching piecewise-linear curve, not measured data.
% (episode,return): (0,.9),(20,.9),(20,.3),(25,.3),(30,.55),
% (35,.75),(40,.85),(50,.9),(60,.9).
% Change=20, detection=25, recovery=40 at predeclared threshold .85.
\begin{tikzpicture}
 \begin{axis}[width=.98\linewidth,height=67mm,
   font=\figurefont\fontsize{8.5}{11}\selectfont,
   axis lines=left,xmin=0,xmax=60,ymin=0,ymax=1.18,
   xtick={0,20,25,40,60},ytick={0,.3,.6,.9},
   xlabel={部署回合$k$},ylabel={假设回报},
   tick label style={font=\figurefont\fontsize{8}{10}\selectfont}]
  \path[fill=BookAmber!12] (axis cs:20,.9) -- (axis cs:40,.9)
    -- (axis cs:40,.85) -- (axis cs:35,.75) -- (axis cs:30,.55)
    -- (axis cs:25,.3) -- (axis cs:20,.3) -- cycle;
  \draw[BookRule,dashed] (axis cs:0,.85) -- (axis cs:60,.85);
  \draw[BookAmber,densely dashed] (axis cs:20,0) -- (axis cs:20,1.05);
  \draw[BookInk,dotted,line width=.8pt] (axis cs:25,0) -- (axis cs:25,.98);
  \draw[BookTeal,dash dot] (axis cs:40,0) -- (axis cs:40,1.05);
  \addplot[BookInk,line width=1pt] coordinates
    {(0,.9)(20,.9)(20,.3)(25,.3)(30,.55)(35,.75)(40,.85)(50,.9)(60,.9)};
  \node[anchor=east] at (axis cs:20,1.09) {变化};
  \node[anchor=west] at (axis cs:25,1.00) {检测};
  \node[anchor=south] at (axis cs:40,1.06) {恢复};
  \node[anchor=west,align=left] at (axis cs:31,.36)
    {检测延迟：5回合\\检测后恢复：15回合};
  \node[anchor=west] at (axis cs:1,.77) {恢复阈值$0.85$};
  \draw[{Latex[length=1.6mm]}-{Latex[length=1.6mm]},BookAmber]
    (axis cs:17,.3) -- (axis cs:17,.9);
  \node[anchor=east,align=right] at (axis cs:16,.48) {性能\\下降};
 \end{axis}
\end{tikzpicture}

 \caption{变化、检测与恢复的教学时间线。曲线由图源中给定的分段线性控制点生成；变化在回合20，检测在25，首次达到预设回报阈值$0.85$在40。下降、延迟和恢复期间的差额各有含义，非任何算法的实测轨迹。}
 \label{fig:rl-transfer-tracking}
\end{figure}

若系统从不按回合重置，前期动作会影响变化发生时的物理状态，就应使用适合持续过程的平均回报、动态比较规则或跟踪损失，不能原样套用共同$\rho_0$的式~\eqref{eq:rl-transfer-dynamic-regret}。本节关注变化期间的决策代价；旧能力保留、遗忘和一般持续学习理论属于第三卷“知识的演进与迁移”部分。

% Outline: 6.5.7
\subsection{任务信息与适应权限的组合边界}
\label{sec:rl-transfer-combinations}

显式描述、记忆、梯度更新与持续交互可以出现在同一系统中。目标坐标已知，只排除了目标身份的不确定性；陌生地面仍可能要求估计摩擦参数。隐藏风向没有描述，也不意味着必须改网络权重：记忆中积累的偏移证据足以改变条件策略。

表~\ref{tab:rl-transfer-permissions}~以信息条件为行、机制权限为列。每格写的是该组合所需的证据和接口，而非给方法贴一个永久类别标签。

\begin{table*}[t]
 \centering\small
 \begin{tabularx}{\linewidth}{@{}>{\raggedright\arraybackslash}p{18mm}*{4}{>{\raggedright\arraybackslash}X}@{}}
  \toprule
  信息条件 & 零样本调用 & 上下文更新 & 参数更新 & 持续交互\\
  \midrule
  显式描述 & 描述语义、范围与接口兼容；不读目标反馈学习 & 描述加获准历史；反馈可补描述遗漏 & 描述加日志或新样本；允许反传 & 允许探测描述未给出的动力学；全部计费\\
  隐藏任务 & 只能基于先验或已有信息选择；不可假装知道任务 & 历史须含区分性反馈；规定窗口与重置 & 更新样本须支持估计；不使用隐藏标签 & 可达的区分动作、反馈权限与恢复预算\\
  \bottomrule
 \end{tabularx}
 \caption{任务信息与适应权限的组合。各列可同时启用，持续交互是一种数据权限，不是与记忆或梯度互斥的网络结构。}
 \label{tab:rl-transfer-permissions}
\end{table*}

例如，显式目标加两步梯度更新可以用目标选择损失，再由新转移修正行动者；但两步所用交互、学习率选择和更新后评价仍须纳入预算。隐藏风向加固定参数记忆则可用第一回合判断风向、第二回合选路；它不消耗梯度步，却消耗反馈和上下文。两种组合只有在总交互及允许信息可比时，才有理由比较“适应更快”。

技能复用也遵循这一划分。任务描述或信念决定高层在合法启动集内选哪项技能，技能结束后根据新观测再组合下一项；若允许更新，可以只改高层选择器，也可以调整技能内部控制参数。后者可能改变持续时间与终止分布，应重新检查与高层价值估计的接口，而不能把内部更新当作无影响替换。技能定义与单任务时间抽象仍沿用第~\ref{chap:rl-learning}~章。

比较方法前应同时固定目标交互数、反馈字段、上下文容量、梯度步数及额外计算。主动辨识与基于测试成绩的选模都使用目标信息，不能从预算中删除。适应机制的组合越灵活，实验协议越需要明确。

% Outline: 6.6
\section{迁移结论的验证与修正}
\label{sec:rl-transfer-validation}

前面给出了哪些知识可能共享、目标任务提供什么信息，以及策略如何形成。验证要将这些条件变成可执行协议：任务从哪里来，哪些反馈在何时可见，结果与哪种比较对象相对。否则一个较高分数无法说明究竟是哪项迁移主张得到了支持。

% Outline: 6.6.1
\subsection{由目标任务适用范围确定任务级数据划分}
\label{sec:rl-transfer-split}

要检验支持内新任务，应把完整任务身份留出；要检验组合泛化，应留出目标与噪声的指定组合；要检验范围外能力，应把相应因素范围从训练和调参任务中排除。若只把同一风向、同一目标下的轨迹随机分成两组，测试的仍是任务内新经历，不是跨任务泛化。

图~\ref{fig:rl-transfer-split}~按任务、回合、片段三层展示这一差别。同一回合的相邻片段共享物理轨迹，同一任务的不同回合共享隐藏参数；把它们都视为独立任务会严重夸大证据量。

\begin{figure}[htbp]
 \centering
 % Original nested sampling diagram; schematic boxes are not sample counts.
% Final width 112 mm. Task split precedes episode and fragment sampling.
\begin{tikzpicture}[x=\linewidth,y=1mm,
 font=\figurefont\fontsize{8.5}{11}\selectfont,
 task/.style={draw=BookRule,line width=.8pt,fill=BookPaper},
 episode/.style={draw=BookRule,line width=.7pt,fill=white}]
 \foreach \x/\name in {.01/训练任务,.35/验证任务,.69/测试任务} {
  \draw[task] (\x,0) rectangle ++(.30,-57);
  \node[anchor=north] at (\x+.15,-2) {\name};
  \foreach \y/\episode in {-13/回合1,-34/回合2} {
   \draw[episode] (\x+.025,\y) rectangle ++(.25,-18);
   \node[anchor=north] at (\x+.15,\y-1) {\episode};
   \foreach \dx/\seg in {.04/片段,.16/片段} {
    \draw[draw=BookTeal,line width=.7pt,fill=BookTeal!10]
      (\x+\dx,\y-10) rectangle ++(.10,-6);
    \node at (\x+\dx+.05,\y-13) {\seg};
   }
  }
 }
 \draw[BookAmber,dashed,line width=.9pt] (.33,4) -- (.33,-60);
 \draw[BookAmber,dashed,line width=.9pt] (.67,4) -- (.67,-60);
 \node[anchor=south,align=center] at (.50,7)
   {在任务层划分：新任务、组合或因素范围};
 \node[anchor=north,align=center,text=BookMuted] at (.50,-63)
   {同回合片段共享轨迹；同任务回合共享任务参数\\
    片段数量不等于独立测试任务数量};
\end{tikzpicture}

 \caption{任务级划分与任务内采样。训练、验证、测试先按所声明的任务关系划分，再在每项任务内采集回合和片段；同色框中片段共享任务参数，不能按片段数声称独立任务数。图为抽样层级示意，不对应某个数据集规模。}
 \label{fig:rl-transfer-split}
\end{figure}

例如，三个独立生成的风向任务各采10回合，每回合切20段，共有600个片段，但任务层独立单位仍只有3个。这些是教学计数。若要评估任务平均回报的不确定性，应在任务层重采样，或用明确的分层统计模型；训练种子另是模型训练随机性的层级，不能将同一个训练种子下的600段当作600次独立训练。

测试协议应预先固定是否允许交互、奖励到达时刻、回合结束与任务切换的记忆重置，以及检查点数量。非平稳测试还要固定任务时间顺序；打乱时间后进行批量拟合，使用的是另一种信息权限。不能让策略提前看到未来风向标签，也不能在恢复时间统计中删除失败回合。

泄漏不只来自重复轨迹。程序化生成器可能只更换纹理而复用同一路线模板，共享随机种子可能重现测试布局，目标描述可能直接包含最优动作编码。应检查生成器结构、种子谱系和描述字段。允许在一项测试任务内按既定协议适应，不等于允许根据整套测试成绩反复重设计训练方法；后一种选择应在验证任务上完成。

% Outline: 6.6.2
\subsection{在任务划分与固定预算下取得性能证据}
\label{sec:rl-transfer-evidence}

初始行为、学习过程和学习完成后的策略回答不同问题。令$c_i$为第$i$个测试任务，$\pi_{i,0}$为未使用该任务更新数据时的策略；$\pi_{i,B}$为按协议使用$B$步反馈后形成的策略。若使用本章前面的无限折扣任务，回报仍为未归一化$J_{c_i}$；若评价固定长度回合，则另注明时域、折扣和截断规则。

\term{零样本回报}为
$n^{-1}\sum_iJ_{c_i}(\pi_{i,0})$。\term{适应后回报}为
$n^{-1}\sum_iJ_{c_i}(\pi_{i,B})$，需要固定适应所得参数与任务信息、在独立评价回合上估计，并另报这部分交互数。评价时仍可按协议更新控制所需的物理状态估计，但不继续辨识任务；若一个记忆器无法分开这两项功能，就应明确评价中仍在适应。\term{预算内适应曲线}则记录随着累计目标交互增加，当前在线回合实际取得的回报，或在指定检查点另测的回报；这两种曲线不能混标，后者的查询成本也不能省略。

为评价学习期间的损失，固定同一目标任务、同一回合时域与初始分布，取指定比较策略$\pi_i^{\text{参照}}$。用$\pi_i^{(k)}$表示第$k$个适应回合的行为规则，区别于使用$B$步反馈后的$\pi_{i,B}$。在$M$个适应回合中定义
\begin{equation}
 \operatorname{AReg}_{i,M}
 =\sum_{k=0}^{M-1}
 \left[J_{c_i}(\pi_i^{\text{参照}})
       -J_{c_i}(\pi_i^{(k)})\right].
 \label{eq:rl-transfer-adaptation-regret}
\end{equation}
若比较者只是一条训练得到的策略，该量可能为负，名称不赋予其最优性；必须写出参照者及信息权限。对于任务本身变化的情形，应另用式~\eqref{eq:rl-transfer-dynamic-regret}~及检测、恢复指标，不能把固定$c_i$的本式直接改个标题。

表~\ref{tab:rl-transfer-baselines}~给出比较协议。$B,K,L,C$分别表示目标适应步、梯度步、上下文长度和额外计算上限；它们是符号预算，不是某篇论文配置。共享训练的前期成本也应报告，即使目标任务比较只匹配部署预算。

\begin{table*}[t]
 \centering\small
 \begin{tabularx}{\linewidth}{@{}>{\raggedright\arraybackslash}p{23mm}*{3}{>{\raggedright\arraybackslash}X}@{}}
  \toprule
  比较对象 & 已有知识与信息 & 目标预算 & 解释范围\\
  \midrule
  从头学习 & 随机初始化，不用源任务知识 & 相同$B,K,L,C$，独立评价另计 & 共享训练是否减少目标学习代价\\
  无适应 & 相同源模型，关闭任务更新 & 适应交互0、梯度0；执行预算相同 & 初始能力与适应带来的变化\\
  单任务训练 & 不共享参数的指定任务训练 & 报预训练量；目标比较仍匹配$B,K,L,C$ & 共享收益与容量、训练量差异\\
  已知任务身份 & 额外给真实任务参数 & 其余预算匹配，信息权限另标 & 知情参考；训练策略未必是上界\\
  可证明最优值 & 完整已知模型或可证明的信息松弛 & 比较者权限显式列出 & 条件成立时才是性能上界\\
  \bottomrule
 \end{tabularx}
 \caption{性能证据需要的参照对象。“无适应”是刻意关闭反馈学习的消融，不能假装它花掉了与适应方法相同的学习预算；单任务参考若用目标数据预训练，必须披露这项额外信息。}
 \label{tab:rl-transfer-baselines}
\end{table*}

图~\ref{fig:rl-transfer-evaluation-curves}~给出一个完全可复算的教学汇总。设目标任务$i=1,\ldots,5$、训练种子$j=1,\ldots,3$，定义假设回报
\[
 \begin{aligned}
 y_{i,j}(B)&=0.2+0.7(1-e^{-B/30})+u_i e^{-B/60}+v_j,\\
 u_i&\in\{-0.16,-0.08,0,0.08,0.16\},\\
 v_j&\in\{-0.02,0,0.02\}.
 \end{aligned}
\]
此处所有数值都是作者设定的回报尺度；每条曲线是给定公式，不是实验测量。先对种子平均，再对任务平均，得到均值$m(B)=0.2+0.7(1-e^{-B/30})$。阴影取种子平均后任务分数的均值加减一个样本标准差，即
$s_{\text{任务}}(B)=\sqrt{0.016}\,e^{-B/60}$，不是均值的置信区间。实际实验还需估计评价回合随机性，并说明任务与训练种子的嵌套或交叉方式。

\begin{figure}[htbp]
 \centering
 % Teaching data, not a benchmark: five tasks crossed with three training seeds.
% y_{i,j}(B)=.2+.7*(1-exp(-B/30))+u_i*exp(-B/60)+v_j.
% u=[-.16,-.08,0,.08,.16], v=[-.02,0,.02].
% After averaging seeds: m=.2+.7*(1-exp(-B/30)), task sample SD=sqrt(.016)*exp(-B/60).
% Points B=[0,10,20,30,40,50,60,80,100,120], rounded to six decimals.
% Line and polygon use linear interpolation between these checkpoints.
\begin{tikzpicture}
 \begin{axis}[width=.98\linewidth,height=63mm,
   font=\figurefont\fontsize{8.5}{11}\selectfont,
   axis lines=left,xmin=0,xmax=120,ymin=0,ymax=1.04,
   xtick={0,30,60,90,120},ytick={0,.2,.4,.6,.8,1},
   xlabel={目标适应交互$B$（步）},ylabel={假设回报},
   tick label style={font=\figurefont\fontsize{8}{10}\selectfont}]
  \addplot[draw=none,fill=BookTeal!16] coordinates
    {(0,.326491)(10,.505500)(20,.631243)(30,.719205)(40,.780425)
     (50,.822760)(60,.851799)(80,.884704)(100,.898919)(120,.904298)
     (120,.870060)(100,.851137)(80,.818019)(60,.758732)(50,.712814)
     (40,.650539)(30,.565764)(20,.449973)(10,.291356)(0,.073509)} \closedcycle;
  \addplot[BookTeal,line width=1pt,mark=*,mark size=1.4pt] coordinates
    {(0,.200000)(10,.398428)(20,.540608)(30,.642484)(40,.715482)
     (50,.767787)(60,.805265)(80,.851362)(100,.875028)(120,.887179)};
  \node[anchor=west,align=left] at (axis cs:5,.92)
    {5项假设任务\\每项3个训练种子};
  \node[anchor=west,align=left] at (axis cs:44,.30)
    {实线：任务均值\\阴影：任务间$\pm1$标准差\\不是均值置信区间};
 \end{axis}
\end{tikzpicture}

 \caption{固定预算下均值与任务间离散程度的教学曲线。横轴为目标适应交互步数，纵轴为假设回报；线为5任务、3训练种子的等权均值，阴影为先平均种子后的任务均值$\pm1$样本标准差。图源按正文公式在所列检查点计算并线性连接，无算法实测排名；真实协议中的评价交互还须单独计入。}
 \label{fig:rl-transfer-evaluation-curves}
\end{figure}

不同任务的回报尺度不同时，可以分别报告原始分数，或使用预先固定、具有解释意义的参照范围归一化；归一化分母接近0的任务需要单独处理。不能观察测试结果后选择有利的缩放方法。汇总还应报告各任务组的失败比例、差异幅度及不确定性，避免少数高分任务掩盖多数任务无收益。训练得到的知情策略即使表现很好，也只有参考意义，不能未经证明称为最优上界。

% Outline: 6.6.3
\subsection{从迁移结果到共享假设的诊断与修正}
\label{sec:rl-transfer-diagnosis}

相对于明确的不共享或从头学习基线，在相同目标信息与预算下，提高指定评价指标称为\term{正迁移}（positive transfer），降低则称为\term{负迁移}（negative transfer）。这里的正负依赖任务组、指标和预算：学习早期有益、最终性能稍低，可以同时成立。不能把一次更高的样本均值直接写成稳定的正迁移结论。

考虑四项等权教学任务，基线回报均为$0.6$，共享模型回报分别为$0.8,0.8,0.8,0.2$。共享均值为$0.65$，平均提高$0.05$，第四项任务却下降$0.4$。若第四项恰好使用未覆盖的左风，这一结果提示检查任务分层与覆盖；它本身尚不能证明动力学共享假设错误。

诊断需要补充可以区分解释的证据。数据覆盖不足时，应控制任务采样和相关状态动作访问；表示容量不足时，应比较参数量和计算量匹配的模型，并检查目标信息是否被丢弃；优化冲突时，可比较分任务梯度的内积、独立训练与任务权重调整。负梯度内积只是局部冲突信号，不自动说明两个任务在真实结构上不兼容。要检验共享结构，可保持其余训练条件，替换被怀疑的共同模块为任务条件或独立模块，并在受影响任务组上重新评价。

图~\ref{fig:rl-transfer-diagnosis}~将修订对象分开。匹配容量后仅改变任务权重便恢复性能，更支持训练机制解释；在充分覆盖与多种优化设置下，共享转移模块仍系统性地把相反风向混合，而任务条件模块消除了相应预测误差，才为重新检查共享假设提供较直接的证据。

\begin{figure}[htbp]
 \centering
 % Original diagnostic graph. Branches are hypotheses, not automatic diagnoses.
% Final width 112 mm; interventions return to knowledge/training decisions.
\begin{tikzpicture}[x=\linewidth,y=1mm,
 font=\figurefont\fontsize{8.5}{11}\selectfont,
 box/.style={draw=BookRule,fill=BookPaper,line width=.8pt,
   minimum height=11mm,align=center,inner sep=1.5mm},
 flow/.style={-{Latex[length=2mm]},draw=BookInk,line width=.9pt}]
 \node[box,text width=40mm] (loss) at (.5,0) {匹配预算后出现性能退化};
 \node[box,text width=61mm] (evidence) at (.5,-24)
   {补充诊断证据\\覆盖、容量、梯度冲突、模块替换};
 \node[box,text width=33mm] (training) at (.20,-53)
   {训练机制失效假设\\采样、表示或优化};
 \node[box,text width=33mm,fill=BookTeal!8] (structure) at (.80,-53)
   {共享结构失效假设\\动力学或接口不兼容};
 \node[box,text width=33mm] (trainfix) at (.20,-80)
   {修订训练方式\\重新做预算匹配比较};
 \node[box,text width=33mm] (structfix) at (.80,-80)
   {重估知识对象与范围\\重新设计任务划分};
 \draw[flow] (loss) -- (evidence);
 \draw[flow] (evidence.south) -- (.5,-38) -| (training.north);
 \draw[flow] (.5,-38) -| (structure.north);
 \draw[flow] (training) -- (trainfix);
 \draw[flow] (structure) -- (structfix);
 \node[anchor=north,text=BookMuted] at (.5,-91)
   {单次失败不足以确定根因；改动后重新验证。};
\end{tikzpicture}

 \caption{由退化结果形成可区分的诊断。分支表示待检验解释，不能只凭一次失败直接确定根因；训练问题修订采样、容量或优化，结构问题则重估共享对象及适用任务。图为本章诊断关系的原创示意。}
 \label{fig:rl-transfer-diagnosis}
\end{figure}

单一基线可能容量不足，一次种子可能恰好失败，模块替换也可能同时增加容量。因此需要逐项控制差异，报告重复试验与不确定性，而不是见到改进就倒推唯一原因。共享结构一旦改变，原来的任务支持与组合假设也要重查，不能沿用旧条件下的泛化声明。

% Outline: 6.6.4
\subsection{由评测证据限定泛化与适应声明}
\label{sec:rl-transfer-claims}

一个可检查的迁移声明至少包含目标任务关系、可用信息、适应预算和性能证据。省去其中任何一项，都可能把更容易的实验解释成更强的能力。表~\ref{tab:rl-transfer-claims}~汇总本章导航变体分别能够支持的说法；表中的结论以相应实验确实取得可靠证据为前提，本章教学数值本身不构成算法实验证据。

\begin{table*}[t]
 \centering\small
 \begin{tabularx}{\linewidth}{@{}>{\raggedright\arraybackslash}p{22mm}*{4}{>{\raggedright\arraybackslash}X}@{}}
  \toprule
  目标任务关系 & 可用信息 & 适应预算 & 评价证据 & 允许的结论\\
  \midrule
  支持内新目标 & 兼容的目标坐标 & 无反馈适应 & 任务级留出零样本回报 & 声明总体内的描述条件泛化\\
  未见目标与风组合 & 相同描述字段 & 固定$B,K,L,C$ & 留出联合组合的分组结果 & 指定因素与组合上的能力\\
  未见噪声范围 & 描述或获准历史 & 明确探测与更新 & 范围外分组、失败与误差 & 已测范围内的外推证据\\
  固定隐藏风向 & 无真实风标签，反馈逐步可见 & 探测、恢复、评价均计 & 辨识误差、适应回报与参照差额 & 规定信息条件下的固定任务适应\\
  变化风向序列 & 不提前见未来变化 & 按时间顺序持续交互 & 动态遗憾、检测延迟与恢复 & 指定变化预算及信号下的跟踪\\
  \bottomrule
 \end{tabularx}
 \caption{证据与声明的对应关系。支持内、组合、范围外、固定隐藏任务和变化任务不是逐级自动推出的能力；每一行都需要自身的任务划分与协议。}
 \label{tab:rl-transfer-claims}
\end{table*}

训练任务数量很大，但每项只改变同一地图的颜色时，测试高分仍不能证明新几何泛化。目标和风向都单独测过，但从未共同改变时，不能推断组合有效。若变化检测器提前读到了下一回合任务标签，快速恢复也不能支持无预告持续适应。这些反例没有否认已测分数，而是限定分数回答的问题。

还应区分平均能力与最低要求。对声明任务分布取得较高平均回报，不等于在所有允许扰动下都不低于某个阈值；后者需要明确扰动集合、风险或最坏情形准则，以及相应验证。鲁棒性与可信保证留到第~\ref{chap:rl-trustworthy}~章讨论，不能由本章经验泛化结果直接推出。

% Outline: 6.7
\section{本章小结}
\label{sec:rl-transfer-summary}

回到章首的小车：换目标时，共同移动规律可能继续使用，价值却要按新奖励重算；地面变滑时，目标含义仍在，旧后果预测却可能失效。判断能否迁移，需要把复用对象与共享条件成对说明，再检查新任务是否属于这些条件覆盖的范围。参数形状相同、平均回报相近或任务编号相邻，都不能替代这一步。

训练阶段形成表示、模型、价值、策略初始化和技能；部署阶段根据描述、历史与权限形成具体策略。固定参数仍可在记忆中适应，参数更新仍需有效数据，持续交互也只有在反馈可识别且动作获准时才补得上信息。后继特征的价值分解依赖共同动力学与奖励特征，GPI比较的是已有策略，贝叶斯任务更新依赖固定任务和正确似然；这些条件说明每项机制解决了哪一部分问题，也留下了哪些工作。

认识固定未知任务与追踪变化任务并不相同。每轮独立变化的两臂例子表明，计算或记忆无法补出缺失的当轮信息。任务级划分、完整预算、匹配比较者和分层失败诊断，共同约束泛化与适应声明。伙伴与对手也会带来任务变化，但各自的目标与共同学习需要联合建模。第~\ref{chap:rl-multi-agent}~章将考察这些依赖，并把任务信息与适应预算的约定用于交互双方。
